% Complete reviewed Unit 10; editorial master: chapters/ch10/student.md
% Markdown SHA-256: 2afdc5fd37c0b5502061641ec5bda6458760b54f69a51c81aa9b1be9ccdc5606
% One-use initial assembly followed by complete content and page review.
\newcommand{\runningtitle}{Unit 10. The case system}
% Layout inherited from the reviewed previous course; set a running title per unit.
% This is a small hand-maintained TeX layout, not a conversion pipeline.
\RequirePackage{fix-cm}
\documentclass[11pt,twoside]{article}
\usepackage{fontspec}
\setmainfont{Cambria}[
  Renderer=HarfBuzz,
  BoldFont=cambriab.ttf,
  ItalicFont=cambriai.ttf,
  BoldItalicFont=cambriaz.ttf
]
\setmonofont{lmmono10-regular.otf}[BoldFont=lmmonolt10-bold.otf]
\usepackage[paperwidth=165mm,paperheight=241mm,inner=21mm,outer=19mm,
  top=22mm,bottom=20mm,headheight=14pt,headsep=7mm,footskip=10mm]{geometry}
\usepackage{xcolor,array,booktabs,longtable,tabularx}
\usepackage{enumitem,fancyhdr}
\usepackage[unicode,hidelinks]{hyperref}
\definecolor{ink}{HTML}{000000}
\definecolor{muted}{HTML}{4C5558}
\setlength{\parindent}{0pt}
\setlength{\parskip}{3.5pt}
\AtBeginDocument{\fontsize{11}{14.3}\selectfont}
\setlength{\tabcolsep}{4pt}
\renewcommand{\arraystretch}{1.10}
\setlist[itemize]{leftmargin=1.3em,label=\textbullet,itemsep=2pt,topsep=3pt,parsep=0pt}
\setlist[enumerate]{leftmargin=1.8em,itemsep=3pt,topsep=4pt,parsep=0pt}
\widowpenalty=10000
\clubpenalty=10000
\emergencystretch=1em
\raggedbottom
\pagestyle{fancy}
\fancyhf{}
\providecommand{\runningtitle}{Old English}
\fancyhead[LE,RO]{\small\itshape \runningtitle}
\fancyhead[LO,RE]{\small Old English}
\fancyfoot[LE,RO]{\small\thepage}
\renewcommand{\headrulewidth}{0.3pt}
\newcommand{\booktitle}[2]{
  {\fontsize{9}{11}\selectfont\color{muted}#1\par}
  \vspace{3pt}
  {\fontsize{20}{24}\selectfont\color{ink}#2\par}
  \vspace{6pt}
}
\newcommand{\source}[1]{\par{\footnotesize\color{muted}#1\par}}
\newcommand{\readingtext}[1]{
  \par\addvspace{5pt}
  \begingroup\leftskip=1.5em
    \fontsize{12}{16}\selectfont #1\par
  \endgroup\addvspace{5pt}
}
\makeatletter
\renewcommand{\section}{\@startsection{section}{1}{\z@}
  {-2.4ex \@plus -0.6ex \@minus -0.2ex}{0.8ex \@plus 0.2ex}
  {\normalfont\fontsize{12}{15}\selectfont\bfseries\color{ink}}}
\renewcommand{\subsection}{\@startsection{subsection}{2}{\z@}
  {-1.8ex \@plus -0.5ex \@minus -0.2ex}{0.5ex \@plus 0.2ex}
  {\normalfont\fontsize{11}{14}\selectfont\bfseries\color{ink}}}
\makeatother

\hyphenpenalty=10000
\exhyphenpenalty=10000
\emergencystretch=3em
\newcommand{\itempara}[2]{\par\addvspace{3pt}\noindent\hangindent=1.8em\hangafter=1\makebox[1.8em][l]{#1.}#2\par}
\newcommand{\reserve}[1]{\par\penalty-100\begingroup\dimen0=\pagegoal\advance\dimen0 by-\pagetotal\ifdim\dimen0<#1\newpage\fi\endgroup}
\begin{document}
\addfontfeatures{Ligatures=TeXOff}
\raggedright
\booktitle{UNIT 10}{The case system in extended sentences}

% mdparagraph:P001:begin
Case identifies a noun phrase's function in its clause and the relation of one noun to another. A dative recipient may contain a genitive possessor; a verb may itself require a genitive object. The same English word “with” may express accompaniment or means. This unit brings these distinctions together in longer sentences.
% mdparagraph:P001:end

% mdparagraph:P002:begin
The fourteen translated examples, the exercises marked F, A, T, R and S, and “The stranger and the monk” are new teaching material. The three Cook exercise sets retain their original English prompts. “The Day of Judgment” is an exact source passage with supplied help and Cook's original notes.
% mdparagraph:P002:end

% mdparagraph:P003:begin
In the tables, 1, 2 and 3 denote person. A repeated spelling occupies a separate grammatical cell in each position. The ordinary nominative, genitive, dative and accusative forms are for use in composition unless expressly marked \textbf{reference}. Reference forms are for recognition; supplied expressions may be used with their printed help. Keep vowel lengths and the distinction between e and ę.
% mdparagraph:P003:end

\reserve{7\baselineskip}
% mdparagraph:P004:begin
\section*{1. The case of a phrase and the dependent genitive}
% mdparagraph:P004:end

% teaching: NEW-U10-GRAM-u10-genitive-relations
% mdparagraph:P005:begin
A noun phrase receives its case from its role in the clause. Its head noun and agreeing modifiers show that case. A genitive dependent inside the phrase keeps the case required by its relation to the head.
% mdparagraph:P005:end

% mdparagraph:P006:begin
A dependent genitive may identify a possessor or a kinship relation, the source or belonging of something, or the person associated with an action named by a noun. In a definite possessive expression, this unit normally puts the genitive before the head: \textbf{ðæs cuman sunu}, “the stranger's son”. The article belongs to cuman; do not add a second article before sunu in this selected pattern.
% mdparagraph:P006:end

% mdparagraph:P007:begin
The head can change case without changing its dependent genitive.
% mdparagraph:P007:end

\reserve{8\baselineskip}
% example: NEW-U10-EXAMPLE01
% mdparagraph:P008:begin
Se munuc sealde ðæs cuman suna męte.
% mdparagraph:P008:end

% mdparagraph:P009:begin
The monk gave food to the stranger’s son.
% mdparagraph:P009:end

% mdparagraph:P010:begin
The whole phrase \textbf{ðæs cuman suna} is the dative recipient of sealde. Its head suna is dative singular, while ðæs cuman is genitive singular, “of the stranger”. Męte, the food given, is accusative. The genitive does not make the complete recipient phrase genitive.
% mdparagraph:P010:end

% mdparagraph:P011:begin
With a noun naming an action or feeling, a dependent genitive may denote the person performing it, a \textbf{subjective genitive}, or the person or thing toward which it is directed, an \textbf{objective genitive}. Determine the relation from the meaning: “the person's fear” can mean fear felt by that person, while “fear of the person” can mean that someone fears that person. The endings alone do not distinguish these readings. See source note 2 for Cook's objective-genitive example and correction C02.
% mdparagraph:P011:end

% teaching: NEW-U10-GRAM-u10-governed-case-contrast
% teaching: NEW-U10-GRAM-u10-forwiernan-government
% teaching: NEW-U10-GRAM-u10-unnan-government
% teaching: NEW-U10-VOC-FORWIERNAN
% teaching: NEW-U10-VOC-UNNAN
% mdparagraph:P012:begin
A \textbf{governed genitive} depends directly on a verb. It is not a possessor inside an object phrase. Keep the earlier constructions: \textbf{gīeman + genitive}, “rule over”, and \textbf{geswīcan + genitive}, “cease from” an action or work. An English object does not by itself justify an Old English accusative.
% mdparagraph:P012:end

% mdparagraph:P013:begin
Two further verbs take both a dative person and a genitive thing:
% mdparagraph:P013:end

% mdparagraph:P014:begin
\textbf{Forwiernan}, “deny, refuse”: the person denied is dative; the thing refused is genitive.
% mdparagraph:P014:end

% mdparagraph:P015:begin
\textbf{Unnan}, “grant, allow”: the person receiving the grant is dative; the thing granted is genitive.
% mdparagraph:P015:end

% mdparagraph:P016:begin
These case frames remain the same with a statement, a question or a command. The complete verb forms follow in §7.
% mdparagraph:P016:end

\reserve{8\baselineskip}
% example: NEW-U10-EXAMPLE02
% mdparagraph:P017:begin
Ðū forwiernst mē mętes.
% mdparagraph:P017:end

% mdparagraph:P018:begin
You deny me food.
% mdparagraph:P018:end

% mdparagraph:P019:begin
Mē is the dative person; mętes is the genitive of the food refused. Compare sęllan, “give”: its person is also dative, but the thing given is accusative. Similar English word order does not make the constructions identical.
% mdparagraph:P019:end

\reserve{7\baselineskip}
% mdparagraph:P020:begin
\section*{2. Nouns: complete models and lexical diagnostics}
% mdparagraph:P020:end

% teaching: NEW-U10-LEXICAL-DIAGNOSTICS
% mdparagraph:P021:begin
Each new noun has its gender, meaning, stem and diagnostic forms below. A \textbf{stem} is the part retained when an ending is added. For a noun assigned to a printed model, use the cell with the same case and number, replace the model stem by the listed stem, and retain that cell's ending. A cell with no ending retains no ending in the new noun. The stated full headword is used wherever the rule specifies it.
% mdparagraph:P021:end

% mdparagraph:P022:begin
The tables give all four cases in both numbers. \textbf{Crīst} is a proper name used only in the singular. \textbf{Dōm}, “judgment”, \textbf{rēaf}, “raiment”, and \textbf{woruld}, “world”, are used here in the singular; their complete plural forms are reference forms, without an additional plural meaning to learn. The remaining nouns have both numbers available in their assigned meanings.
% mdparagraph:P022:end

\reserve{7\baselineskip}
% mdparagraph:P023:begin
\subsection*{Masculine models}
% mdparagraph:P023:end

% mdparagraph:P024:begin
The ordinary consonant masculine uses -es and -e in the singular genitive and dative, and -as, -a, -um, -as in the plural. Its nominative and accusative singular have no ending. \textbf{Dōm} supplies this complete model; its own plural is reference.
% mdparagraph:P024:end

% teaching: NEW-U10-NOM-DOM
% table: NEW-U10-NOM-DOM
% mdtable:T01:begin
\par\begingroup
\fontsize{10.5}{13.2}\selectfont
\setlength{\LTleft}{0pt}\setlength{\LTright}{\fill}
\noindent
\begin{tabular}{@{}>{\raggedright\arraybackslash}p{39.46mm}>{\raggedright\arraybackslash}p{39.46mm}>{\raggedright\arraybackslash}p{39.46mm}@{}}
\toprule
\textbf{Case} & \textbf{Singular} & \textbf{Plural} \\
\midrule
Nominative & dōm & dōmas \\
Genitive & dōmes & dōma \\
Dative & dōme & dōmum \\
Accusative & dōm & dōmas \\
\bottomrule
\end{tabular}
\par\endgroup

% mdtable:T01:end
% mdparagraph:P025:begin
\textbf{Ęngel}, “angel”, retains its internal e in the nominative and accusative singular. Before an ending it has the shorter stem \textbf{ęngl-}. Use this full-headword/\allowbreak{}oblique-stem distinction also for \textbf{dryhten}, stem \textbf{dryhtn-}, and \textbf{dēofol}, stem \textbf{dēofl-}.
% mdparagraph:P025:end

% teaching: NEW-U10-NOM-ENGEL
% table: NEW-U10-NOM-ENGEL
% mdtable:T02:begin
\par\begingroup
\fontsize{10.5}{13.2}\selectfont
\setlength{\LTleft}{0pt}\setlength{\LTright}{\fill}
\noindent
\begin{tabular}{@{}>{\raggedright\arraybackslash}p{39.46mm}>{\raggedright\arraybackslash}p{39.46mm}>{\raggedright\arraybackslash}p{39.46mm}@{}}
\toprule
\textbf{Case} & \textbf{Singular} & \textbf{Plural} \\
\midrule
Nominative & ęngel & ęnglas \\
Genitive & ęngles & ęngla \\
Dative & ęngle & ęnglum \\
Accusative & ęngel & ęnglas \\
\bottomrule
\end{tabular}
\par\endgroup

% mdtable:T02:end
% mdparagraph:P026:begin
Thus every ending-bearing cell of dryhten and dēofol uses the same ending as the corresponding cell of ęngel, attached to dryhtn- or dēofl-. Their nominative and accusative singular are the full dryhten and dēofol. Retained-vowel variants of dryhten, including genitive \textbf{dryhtenes}, are also valid; the selected table route uses dryhtn- consistently. Do not extend this vowel loss to an unlisted noun.
% mdparagraph:P026:end

% mdparagraph:P027:begin
\textbf{Scēaphierde}, “shepherd”, keeps -e in the nominative and accusative singular. Its stem for endings is \textbf{scēaphierd-}; the genitive singular and the plural show the ordinary masculine endings.
% mdparagraph:P027:end

% teaching: NEW-U10-NOM-SCEAPHIERDE
% table: NEW-U10-NOM-SCEAPHIERDE
% mdtable:T03:begin
\par\begingroup
\fontsize{10.5}{13.2}\selectfont
\setlength{\LTleft}{0pt}\setlength{\LTright}{\fill}
\noindent
\begin{tabular}{@{}>{\raggedright\arraybackslash}p{34.82mm}>{\raggedright\arraybackslash}p{41.78mm}>{\raggedright\arraybackslash}p{41.78mm}@{}}
\toprule
\textbf{Case} & \textbf{Singular} & \textbf{Plural} \\
\midrule
Nominative & scēaphierde & scēaphierdas \\
Genitive & scēaphierdes & scēaphierda \\
Dative & scēaphierde & scēaphierdum \\
Accusative & scēaphierde & scēaphierdas \\
\bottomrule
\end{tabular}
\par\endgroup

% mdtable:T03:end
% mdparagraph:P028:begin
The weak masculine \textbf{cuma}, “stranger, guest”, has -a only in the nominative singular, -an in its other singular cases and in the nominative and accusative plural, -ena in the genitive plural, and -um in the dative plural. \textbf{Nama}, stem \textbf{nam-}, follows every cell of this model.
% mdparagraph:P028:end

% teaching: NEW-U10-NOM-CUMA
% table: NEW-U10-NOM-CUMA
% mdtable:T04:begin
\par\begingroup
\fontsize{10.5}{13.2}\selectfont
\setlength{\LTleft}{0pt}\setlength{\LTright}{\fill}
\noindent
\begin{tabular}{@{}>{\raggedright\arraybackslash}p{39.46mm}>{\raggedright\arraybackslash}p{39.46mm}>{\raggedright\arraybackslash}p{39.46mm}@{}}
\toprule
\textbf{Case} & \textbf{Singular} & \textbf{Plural} \\
\midrule
Nominative & cuma & cuman \\
Genitive & cuman & cumena \\
Dative & cuman & cumum \\
Accusative & cuman & cuman \\
\bottomrule
\end{tabular}
\par\endgroup

% mdtable:T04:end
% mdparagraph:P029:begin
\textbf{Fæder}, “father”, has an exceptional endingless singular. Its plural is fæder- plus -as, -a, -um, -as. Genitive \textbf{fæderes} and dative \textbf{fædere} are accepted variants. The selected singular below has fæder in all four cases; do not derive a changed-vowel dative by analogy with brōðor.
% mdparagraph:P029:end

% teaching: NEW-U10-NOM-FAEDER
% table: NEW-U10-NOM-FAEDER
% mdtable:T05:begin
\par\begingroup
\fontsize{10.5}{13.2}\selectfont
\setlength{\LTleft}{0pt}\setlength{\LTright}{\fill}
\noindent
\begin{tabular}{@{}>{\raggedright\arraybackslash}p{39.46mm}>{\raggedright\arraybackslash}p{39.46mm}>{\raggedright\arraybackslash}p{39.46mm}@{}}
\toprule
\textbf{Case} & \textbf{Singular} & \textbf{Plural} \\
\midrule
Nominative & fæder & fæderas \\
Genitive & fæder & fædera \\
Dative & fæder & fæderum \\
Accusative & fæder & fæderas \\
\bottomrule
\end{tabular}
\par\endgroup

% mdtable:T05:end
\reserve{7\baselineskip}
% mdparagraph:P030:begin
\subsection*{Neuter model}
% mdparagraph:P030:end

% mdparagraph:P031:begin
The heavy neuter \textbf{scēap}, “sheep”, has identical nominative and accusative forms in both numbers. The genitive and dative singular have -es and -e; the genitive and dative plural have -a and -um. Context and agreement distinguish singular from plural where the noun itself does not.
% mdparagraph:P031:end

% teaching: NEW-U10-NOM-SCEAP
% table: NEW-U10-NOM-SCEAP
% mdtable:T06:begin
\par\begingroup
\fontsize{10.5}{13.2}\selectfont
\setlength{\LTleft}{0pt}\setlength{\LTright}{\fill}
\noindent
\begin{tabular}{@{}>{\raggedright\arraybackslash}p{39.46mm}>{\raggedright\arraybackslash}p{39.46mm}>{\raggedright\arraybackslash}p{39.46mm}@{}}
\toprule
\textbf{Case} & \textbf{Singular} & \textbf{Plural} \\
\midrule
Nominative & scēap & scēap \\
Genitive & scēapes & scēapa \\
Dative & scēape & scēapum \\
Accusative & scēap & scēap \\
\bottomrule
\end{tabular}
\par\endgroup

% mdtable:T06:end
% mdparagraph:P032:begin
Replace scēap- by \textbf{rēaf-}, \textbf{cweartern-}, \textbf{giesthūs-} or \textbf{setl-} in every cell, retaining the same ending, including the zero ending. Cweartern keeps its final -ern throughout: cwearternes, not an automatically shortened stem. Giesthūs keeps the heavy neuter pattern of its final member hūs. No internal vowel is removed from these stems.
% mdparagraph:P032:end

\reserve{7\baselineskip}
% mdparagraph:P033:begin
\subsection*{Feminine models}
% mdparagraph:P033:end

% mdparagraph:P034:begin
\textbf{Hand}, “hand; side”, has an endingless accusative singular in the selected pattern. Its genitive and dative singular have -e; its nominative, genitive and accusative plural have -a, and the dative plural -um. \textbf{Woruld}, stem woruld-, follows these same eight cells.
% mdparagraph:P034:end

% teaching: NEW-U10-NOM-HAND
% table: NEW-U10-NOM-HAND
% mdtable:T07:begin
\par\begingroup
\fontsize{10.5}{13.2}\selectfont
\setlength{\LTleft}{0pt}\setlength{\LTright}{\fill}
\noindent
\begin{tabular}{@{}>{\raggedright\arraybackslash}p{39.46mm}>{\raggedright\arraybackslash}p{39.46mm}>{\raggedright\arraybackslash}p{39.46mm}@{}}
\toprule
\textbf{Case} & \textbf{Singular} & \textbf{Plural} \\
\midrule
Nominative & hand & handa \\
Genitive & hande & handa \\
Dative & hande & handum \\
Accusative & hand & handa \\
\bottomrule
\end{tabular}
\par\endgroup

% mdtable:T07:end
% mdparagraph:P035:begin
For hand, genitive and dative singular \textbf{handa}, accusative singular \textbf{hande} or \textbf{handa}, and nominative/\allowbreak{}accusative plural \textbf{hande} are accepted variants. Woruld has corresponding variants \textbf{worulda} in the genitive/\allowbreak{}dative singular, \textbf{worulde} or \textbf{worulda} in the accusative singular, and \textbf{worulde} in the nominative/\allowbreak{}accusative plural. Learn the selected pattern first; a variant is not an additional case.
% mdparagraph:P035:end

% mdparagraph:P036:begin
\textbf{Ðēod}, “people, nation”, has -e in the accusative singular as well as the genitive and dative. Its plural otherwise has the same endings as hand. \textbf{Wǣd}, “garment”, replaces ðēod- by wǣd- in all eight cells. Both nouns also allow -e in the nominative/\allowbreak{}accusative plural: ðēode, wǣde.
% mdparagraph:P036:end

% teaching: NEW-U10-NOM-THEOD
% table: NEW-U10-NOM-THEOD
% mdtable:T08:begin
\par\begingroup
\fontsize{10.5}{13.2}\selectfont
\setlength{\LTleft}{0pt}\setlength{\LTright}{\fill}
\noindent
\begin{tabular}{@{}>{\raggedright\arraybackslash}p{39.46mm}>{\raggedright\arraybackslash}p{39.46mm}>{\raggedright\arraybackslash}p{39.46mm}@{}}
\toprule
\textbf{Case} & \textbf{Singular} & \textbf{Plural} \\
\midrule
Nominative & ðēod & ðēoda \\
Genitive & ðēode & ðēoda \\
Dative & ðēode & ðēodum \\
Accusative & ðēode & ðēoda \\
\bottomrule
\end{tabular}
\par\endgroup

% mdtable:T08:end
% mdparagraph:P037:begin
The root noun \textbf{gāt}, “goat”, changes ā to ǣ in the selected dative singular and nominative/\allowbreak{}accusative plural. Its genitive singular has -e; the genitive and dative plural retain ā and add -a and -um. Genitive singular \textbf{gǣt} and dative singular \textbf{gāte} are also valid. The table selects one form per cell; do not make a masculine -as plural or extend ǣ into gātum.
% mdparagraph:P037:end

% teaching: NEW-U10-NOM-GAT
% table: NEW-U10-NOM-GAT
% mdtable:T09:begin
\par\begingroup
\fontsize{10.5}{13.2}\selectfont
\setlength{\LTleft}{0pt}\setlength{\LTright}{\fill}
\noindent
\begin{tabular}{@{}>{\raggedright\arraybackslash}p{39.46mm}>{\raggedright\arraybackslash}p{39.46mm}>{\raggedright\arraybackslash}p{39.46mm}@{}}
\toprule
\textbf{Case} & \textbf{Singular} & \textbf{Plural} \\
\midrule
Nominative & gāt & gǣt \\
Genitive & gāte & gāta \\
Dative & gǣt & gātum \\
Accusative & gāt & gǣt \\
\bottomrule
\end{tabular}
\par\endgroup

% mdtable:T09:end
\reserve{7\baselineskip}
% mdparagraph:P038:begin
\subsection*{Noun vocabulary and model assignments}
% mdparagraph:P038:end

% mdparagraph:P039:begin
In the singular diagnostic column the order is \textbf{genitive /\allowbreak{} dative /\allowbreak{} accusative}. The headword is nominative singular. “Stem; model” specifies the complete table to use above; the nominative plural checks the number pattern. For Crīst use only the singular half of dōm. For dryhten and dēofol use the full headword in nominative/\allowbreak{}accusative singular and the shorter listed stem in every other cell. In every other row, replace the named model's stem directly by the listed stem; the exceptional nouns have their own literal tables.
% mdparagraph:P039:end

% teaching: NEW-U10-VOC-CRIST
% teaching: NEW-U10-NOM-CRIST
% diagnostic: NEW-U10-DIAG-CRIST
% teaching: NEW-U10-VOC-DRYHTEN
% teaching: NEW-U10-NOM-DRYHTEN
% diagnostic: NEW-U10-DIAG-DRYHTEN
% teaching: NEW-U10-VOC-DOM
% diagnostic: NEW-U10-DIAG-DOM
% teaching: NEW-U10-VOC-REAF
% teaching: NEW-U10-NOM-REAF
% diagnostic: NEW-U10-DIAG-REAF
% teaching: NEW-U10-VOC-CUMA
% diagnostic: NEW-U10-DIAG-CUMA
% teaching: NEW-U10-VOC-CWEARTERN
% teaching: NEW-U10-NOM-CWEARTERN
% diagnostic: NEW-U10-DIAG-CWEARTERN
% teaching: NEW-U10-VOC-GIESTHUS
% teaching: NEW-U10-NOM-GIESTHUS
% diagnostic: NEW-U10-DIAG-GIESTHUS
% teaching: NEW-U10-VOC-FAEDER
% diagnostic: NEW-U10-DIAG-FAEDER
% teaching: NEW-U10-VOC-WORULD
% teaching: NEW-U10-NOM-WORULD
% diagnostic: NEW-U10-DIAG-WORULD
% teaching: NEW-U10-VOC-SCEAPHIERDE
% diagnostic: NEW-U10-DIAG-SCEAPHIERDE
% teaching: NEW-U10-VOC-SCEAP
% diagnostic: NEW-U10-DIAG-SCEAP
% teaching: NEW-U10-VOC-SETL
% teaching: NEW-U10-NOM-SETL
% diagnostic: NEW-U10-DIAG-SETL
% teaching: NEW-U10-VOC-HAND
% diagnostic: NEW-U10-DIAG-HAND
% teaching: NEW-U10-VOC-THEOD
% diagnostic: NEW-U10-DIAG-THEOD
% teaching: NEW-U10-VOC-DEOFOL
% teaching: NEW-U10-NOM-DEOFOL
% diagnostic: NEW-U10-DIAG-DEOFOL
% teaching: NEW-U10-VOC-WAED
% teaching: NEW-U10-NOM-WAED
% diagnostic: NEW-U10-DIAG-WAED
% teaching: NEW-U10-VOC-GAT
% diagnostic: NEW-U10-DIAG-GAT
% teaching: NEW-U10-VOC-NAMA
% teaching: NEW-U10-NOM-NAMA
% diagnostic: NEW-U10-DIAG-NAMA
% teaching: NEW-U10-VOC-ENGEL
% diagnostic: NEW-U10-DIAG-ENGEL
% mdtable:T10:begin
\par\begingroup
\fontsize{10.5}{13.2}\selectfont
\setlength{\LTleft}{0pt}\setlength{\LTright}{\fill}
\begin{longtable}{@{}>{\raggedright\arraybackslash}p{37.19mm}>{\raggedright\arraybackslash}p{22.35mm}>{\raggedright\arraybackslash}p{33.52mm}>{\raggedright\arraybackslash}p{22.5mm}@{}}
\toprule
\textbf{Headword and meaning} & \textbf{Stem; model} & \textbf{Singular G /\allowbreak{} D /\allowbreak{} A} & \textbf{Nominative plural} \\
\midrule
\endfirsthead
\toprule
\textbf{Headword and meaning} & \textbf{Stem; model} & \textbf{Singular G /\allowbreak{} D /\allowbreak{} A} & \textbf{Nominative plural} \\
\midrule
\endhead
\midrule
\multicolumn{4}{r}{\footnotesize Continued on the following page.}\\
\endfoot
\bottomrule
\endlastfoot
Crīst (m.), Christ & Crīst; dōm & Crīstes /\allowbreak{} Crīste /\allowbreak{} Crīst & — \\
dryhten (m.), lord; the Lord & dryhtn; ęngel & dryhtnes /\allowbreak{} dryhtne /\allowbreak{} dryhten & dryhtnas \\
dōm (m.), judgment; dōmes dæg, Day of Judgment & dōm; dōm & dōmes /\allowbreak{} dōme /\allowbreak{} dōm & dōmas \\
rēaf (n.), raiment; clothing & rēaf; scēap & rēafes /\allowbreak{} rēafe /\allowbreak{} rēaf & rēaf \\
cuma (m.), stranger; guest & cum; cuma & cuman /\allowbreak{} cuman /\allowbreak{} cuman & cuman \\
cweartern (n.), prison & cweartern; scēap & cwearternes /\allowbreak{} cwearterne /\allowbreak{} cweartern & cweartern \\
giesthūs (n.), inn; guest-house & giesthūs; scēap & giesthūses /\allowbreak{} giesthūse /\allowbreak{} giesthūs & giesthūs \\
fæder (m.), father; also the spiritual relationship & fæder; fæder & fæder /\allowbreak{} fæder /\allowbreak{} fæder & fæderas \\
woruld (f.), world & woruld; hand & worulde /\allowbreak{} worulde /\allowbreak{} woruld & worulda \\
scēaphierde (m.), shepherd & scēaphierd; scēaphierde & scēaphierdes /\allowbreak{} scēaphierde /\allowbreak{} scēaphierde & scēaphierdas \\
scēap (n.), sheep & scēap; scēap & scēapes /\allowbreak{} scēape /\allowbreak{} scēap & scēap \\
setl (n.), seat; throne & setl; scēap & setles /\allowbreak{} setle /\allowbreak{} setl & setl \\
hand (f.), hand; right or left side & hand; hand & hande /\allowbreak{} hande /\allowbreak{} hand & handa \\
ðēod (f.), people; nation & ðēod; ðēod & ðēode /\allowbreak{} ðēode /\allowbreak{} ðēode & ðēoda \\
dēofol (m.), devil; demon & dēofl; ęngel & dēofles /\allowbreak{} dēofle /\allowbreak{} dēofol & dēoflas \\
wǣd (f.), garment & wǣd; ðēod & wǣde /\allowbreak{} wǣde /\allowbreak{} wǣde & wǣda \\
gāt (f.), goat & gāt; gāt & gāte /\allowbreak{} gǣt /\allowbreak{} gāt & gǣt \\
nama (m.), name & nam; cuma & naman /\allowbreak{} naman /\allowbreak{} naman & naman \\
ęngel (m.), angel & ęngl; ęngel & ęngles /\allowbreak{} ęngle /\allowbreak{} ęngel & ęnglas \\
\end{longtable}
\par\endgroup

% mdtable:T10:end
% mdparagraph:P040:begin
The model assignment gives every remaining form. For example, dēofol takes the ending -um of ęnglum after dēofl-, producing \textbf{dēoflum}; the nominative singular remains \textbf{dēofol}. Nama replaces cum- in cumena with nam-, producing \textbf{namena}. Wǣd replaces ðēod- in ðēode with wǣd-, producing \textbf{wǣde} in each corresponding singular oblique cell. Apply the same operation separately to each case and number.
% mdparagraph:P040:end

\reserve{7\baselineskip}
% mdparagraph:P041:begin
\section*{3. Adjectives: agreement, complete patterns and people}
% mdparagraph:P041:end

% mdparagraph:P042:begin
Adjectives agree in gender, number and case with the noun they modify. Use the strong form in the ordinary indefinite and predicate uses already learned; use the weak form with the definite article in a definite attributive phrase. A separately predicated adjective remains strong even when its noun phrase contains an article.
% mdparagraph:P042:end

% mdparagraph:P043:begin
The following tables give complete strong and weak systems. Nominative, genitive, dative and accusative forms are for composition in these uses. The distinct masculine and neuter singular instrumental cells in these new adjective tables are \textbf{reference only}. “As dative” means the form is identical to the dative in the same gender and number. Productive instrumental practice in this unit uses the separately taught gōd and demonstrative forms in §6.
% mdparagraph:P043:end

\reserve{7\baselineskip}
% mdparagraph:P044:begin
\subsection*{The hungrig pattern}
% mdparagraph:P044:end

% mdparagraph:P045:begin
\textbf{Hungrig}, “hungry”, keeps its internal vowel before every ending in the selected forms. In the strong pattern the feminine nominative singular and neuter nominative/\allowbreak{}accusative plural have -u. The masculine accusative singular has -ne; feminine genitive/\allowbreak{}dative singular -re; genitive plural -ra; dative masculine/\allowbreak{}neuter singular and all plural genders -um. The weak pattern has nominative singular -a in the masculine and -e in the other genders, and -an in the other singular cells except neuter accusative -e. Its plural is -an, -ra, -um, -an.
% mdparagraph:P045:end

% teaching: NEW-U10-ADJ-HUNGRIG-STRONG
% table: NEW-U10-ADJ-HUNGRIG-STRONG
% mdtable:T11:begin
% table-reflow: singular-plural-panels
\par\begingroup
\fontsize{10.5}{13.2}\selectfont
\begin{minipage}{\linewidth}\raggedright
\setlength{\LTleft}{0pt}\setlength{\LTright}{\fill}
\noindent
\begin{tabular}{@{}>{\raggedright\arraybackslash}p{22mm}>{\raggedright\arraybackslash}p{31.18mm}>{\raggedright\arraybackslash}p{31.18mm}>{\raggedright\arraybackslash}p{31.18mm}@{}}
\toprule
\textbf{Case} & \textbf{Masculine singular} & \textbf{Feminine singular} & \textbf{Neuter singular} \\
\midrule
Nominative & hungrig & hungrigu & hungrig \\
Genitive & hungriges & hungrigre & hungriges \\
Dative & hungrigum & hungrigre & hungrigum \\
Accusative & hungrigne & hungrige & hungrig \\
Instrumental & hungrige & as dative & hungrige \\
\bottomrule
\end{tabular}
\par\addvspace{5pt}
\setlength{\LTleft}{0pt}\setlength{\LTright}{\fill}
\noindent
\begin{tabular}{@{}>{\raggedright\arraybackslash}p{22mm}>{\raggedright\arraybackslash}p{31.18mm}>{\raggedright\arraybackslash}p{31.18mm}>{\raggedright\arraybackslash}p{31.18mm}@{}}
\toprule
\textbf{Case} & \textbf{Masculine plural} & \textbf{Feminine plural} & \textbf{Neuter plural} \\
\midrule
Nominative & hungrige & hungriga & hungrigu \\
Genitive & hungrigra & hungrigra & hungrigra \\
Dative & hungrigum & hungrigum & hungrigum \\
Accusative & hungrige & hungriga & hungrigu \\
Instrumental & as dative & as dative & as dative \\
\bottomrule
\end{tabular}
\par\addvspace{5pt}
\end{minipage}\par\endgroup
% mdtable:T11:end

% teaching: NEW-U10-ADJ-HUNGRIG-WEAK
% table: NEW-U10-ADJ-HUNGRIG-WEAK
% mdtable:T12:begin
% table-reflow: singular-plural-panels
\par\begingroup
\fontsize{10.5}{13.2}\selectfont
\begin{minipage}{\linewidth}\raggedright
\setlength{\LTleft}{0pt}\setlength{\LTright}{\fill}
\noindent
\begin{tabular}{@{}>{\raggedright\arraybackslash}p{22mm}>{\raggedright\arraybackslash}p{31.18mm}>{\raggedright\arraybackslash}p{31.18mm}>{\raggedright\arraybackslash}p{31.18mm}@{}}
\toprule
\textbf{Case} & \textbf{Masculine singular} & \textbf{Feminine singular} & \textbf{Neuter singular} \\
\midrule
Nominative & hungriga & hungrige & hungrige \\
Genitive & hungrigan & hungrigan & hungrigan \\
Dative & hungrigan & hungrigan & hungrigan \\
Accusative & hungrigan & hungrigan & hungrige \\
Instrumental & hungrigan & as dative & hungrigan \\
\bottomrule
\end{tabular}
\par\addvspace{5pt}
\setlength{\LTleft}{0pt}\setlength{\LTright}{\fill}
\noindent
\begin{tabular}{@{}>{\raggedright\arraybackslash}p{22mm}>{\raggedright\arraybackslash}p{31.18mm}>{\raggedright\arraybackslash}p{31.18mm}>{\raggedright\arraybackslash}p{31.18mm}@{}}
\toprule
\textbf{Case} & \textbf{Masculine plural} & \textbf{Feminine plural} & \textbf{Neuter plural} \\
\midrule
Nominative & hungrigan & hungrigan & hungrigan \\
Genitive & hungrigra & hungrigra & hungrigra \\
Dative & hungrigum & hungrigum & hungrigum \\
Accusative & hungrigan & hungrigan & hungrigan \\
Instrumental & as dative & as dative & as dative \\
\bottomrule
\end{tabular}
\par\addvspace{5pt}
\end{minipage}\par\endgroup
% mdtable:T12:end

\reserve{7\baselineskip}
% mdparagraph:P046:begin
\subsection*{Adjectives with a lexical -e}
% mdparagraph:P046:end

% mdparagraph:P047:begin
\textbf{Dēorwierðe}, “precious, valuable”, has the stem \textbf{dēorwierð-}. Its lexical -e remains in strong masculine/\allowbreak{}neuter nominative singular and neuter accusative singular. Replace that -e before another ending; do not add endings to dēorwierðe-. In the other cells the selected pattern has the endings printed below, including -u in feminine nominative singular and neuter nominative/\allowbreak{}accusative plural.
% mdparagraph:P047:end

% teaching: NEW-U10-ADJ-DEORWIERTH-STRONG
% table: NEW-U10-ADJ-DEORWIERTH-STRONG
% mdtable:T13:begin
% table-reflow: singular-plural-panels
\par\begingroup
\fontsize{10.5}{13.2}\selectfont
\begin{minipage}{\linewidth}\raggedright
\setlength{\LTleft}{0pt}\setlength{\LTright}{\fill}
\noindent
\begin{tabular}{@{}>{\raggedright\arraybackslash}p{22mm}>{\raggedright\arraybackslash}p{31.18mm}>{\raggedright\arraybackslash}p{31.18mm}>{\raggedright\arraybackslash}p{31.18mm}@{}}
\toprule
\textbf{Case} & \textbf{Masculine singular} & \textbf{Feminine singular} & \textbf{Neuter singular} \\
\midrule
Nominative & dēorwierðe & dēorwierðu & dēorwierðe \\
Genitive & dēorwierðes & dēorwierðre & dēorwierðes \\
Dative & dēorwierðum & dēorwierðre & dēorwierðum \\
Accusative & dēorwierðne & dēorwierðe & dēorwierðe \\
Instrumental & dēorwierðe & as dative & dēorwierðe \\
\bottomrule
\end{tabular}
\par\addvspace{5pt}
\setlength{\LTleft}{0pt}\setlength{\LTright}{\fill}
\noindent
\begin{tabular}{@{}>{\raggedright\arraybackslash}p{22mm}>{\raggedright\arraybackslash}p{31.18mm}>{\raggedright\arraybackslash}p{31.18mm}>{\raggedright\arraybackslash}p{31.18mm}@{}}
\toprule
\textbf{Case} & \textbf{Masculine plural} & \textbf{Feminine plural} & \textbf{Neuter plural} \\
\midrule
Nominative & dēorwierðe & dēorwierða & dēorwierðu \\
Genitive & dēorwierðra & dēorwierðra & dēorwierðra \\
Dative & dēorwierðum & dēorwierðum & dēorwierðum \\
Accusative & dēorwierðe & dēorwierða & dēorwierðu \\
Instrumental & as dative & as dative & as dative \\
\bottomrule
\end{tabular}
\par\addvspace{5pt}
\end{minipage}\par\endgroup
% mdtable:T13:end

% teaching: NEW-U10-ADJ-DEORWIERTH-WEAK
% table: NEW-U10-ADJ-DEORWIERTH-WEAK
% mdtable:T14:begin
% table-reflow: singular-plural-panels
\par\begingroup
\fontsize{10.5}{13.2}\selectfont
\begin{minipage}{\linewidth}\raggedright
\setlength{\LTleft}{0pt}\setlength{\LTright}{\fill}
\noindent
\begin{tabular}{@{}>{\raggedright\arraybackslash}p{22mm}>{\raggedright\arraybackslash}p{31.18mm}>{\raggedright\arraybackslash}p{31.18mm}>{\raggedright\arraybackslash}p{31.18mm}@{}}
\toprule
\textbf{Case} & \textbf{Masculine singular} & \textbf{Feminine singular} & \textbf{Neuter singular} \\
\midrule
Nominative & dēorwierða & dēorwierðe & dēorwierðe \\
Genitive & dēorwierðan & dēorwierðan & dēorwierðan \\
Dative & dēorwierðan & dēorwierðan & dēorwierðan \\
Accusative & dēorwierðan & dēorwierðan & dēorwierðe \\
Instrumental & dēorwierðan & as dative & dēorwierðan \\
\bottomrule
\end{tabular}
\par\addvspace{5pt}
\setlength{\LTleft}{0pt}\setlength{\LTright}{\fill}
\noindent
\begin{tabular}{@{}>{\raggedright\arraybackslash}p{22mm}>{\raggedright\arraybackslash}p{31.18mm}>{\raggedright\arraybackslash}p{31.18mm}>{\raggedright\arraybackslash}p{31.18mm}@{}}
\toprule
\textbf{Case} & \textbf{Masculine plural} & \textbf{Feminine plural} & \textbf{Neuter plural} \\
\midrule
Nominative & dēorwierðan & dēorwierðan & dēorwierðan \\
Genitive & dēorwierðra & dēorwierðra & dēorwierðra \\
Dative & dēorwierðum & dēorwierðum & dēorwierðum \\
Accusative & dēorwierðan & dēorwierðan & dēorwierðan \\
Instrumental & as dative & as dative & as dative \\
\bottomrule
\end{tabular}
\par\addvspace{5pt}
\end{minipage}\par\endgroup
% mdtable:T14:end

\reserve{7\baselineskip}
% mdparagraph:P048:begin
\subsection*{Long stems and doubled consonants}
% mdparagraph:P048:end

% mdparagraph:P049:begin
\textbf{Unrihtwīs}, “unrighteous”, has a long final stem syllable. Its strong feminine nominative singular and neuter nominative/\allowbreak{}accusative plural have no -u. Keep the long ī throughout. The other strong and weak endings are shown in full.
% mdparagraph:P049:end

% teaching: NEW-U10-ADJ-UNRIHTWIS-STRONG
% table: NEW-U10-ADJ-UNRIHTWIS-STRONG
% mdtable:T15:begin
% table-reflow: singular-plural-panels
\par\begingroup
\fontsize{10.5}{13.2}\selectfont
\begin{minipage}{\linewidth}\raggedright
\setlength{\LTleft}{0pt}\setlength{\LTright}{\fill}
\noindent
\begin{tabular}{@{}>{\raggedright\arraybackslash}p{22mm}>{\raggedright\arraybackslash}p{31.18mm}>{\raggedright\arraybackslash}p{31.18mm}>{\raggedright\arraybackslash}p{31.18mm}@{}}
\toprule
\textbf{Case} & \textbf{Masculine singular} & \textbf{Feminine singular} & \textbf{Neuter singular} \\
\midrule
Nominative & unrihtwīs & unrihtwīs & unrihtwīs \\
Genitive & unrihtwīses & unrihtwīsre & unrihtwīses \\
Dative & unrihtwīsum & unrihtwīsre & unrihtwīsum \\
Accusative & unrihtwīsne & unrihtwīse & unrihtwīs \\
Instrumental & unrihtwīse & as dative & unrihtwīse \\
\bottomrule
\end{tabular}
\par\addvspace{5pt}
\setlength{\LTleft}{0pt}\setlength{\LTright}{\fill}
\noindent
\begin{tabular}{@{}>{\raggedright\arraybackslash}p{22mm}>{\raggedright\arraybackslash}p{31.18mm}>{\raggedright\arraybackslash}p{31.18mm}>{\raggedright\arraybackslash}p{31.18mm}@{}}
\toprule
\textbf{Case} & \textbf{Masculine plural} & \textbf{Feminine plural} & \textbf{Neuter plural} \\
\midrule
Nominative & unrihtwīse & unrihtwīsa & unrihtwīs \\
Genitive & unrihtwīsra & unrihtwīsra & unrihtwīsra \\
Dative & unrihtwīsum & unrihtwīsum & unrihtwīsum \\
Accusative & unrihtwīse & unrihtwīsa & unrihtwīs \\
Instrumental & as dative & as dative & as dative \\
\bottomrule
\end{tabular}
\par\addvspace{5pt}
\end{minipage}\par\endgroup
% mdtable:T15:end

% teaching: NEW-U10-ADJ-UNRIHTWIS-WEAK
% table: NEW-U10-ADJ-UNRIHTWIS-WEAK
% mdtable:T16:begin
% table-reflow: singular-plural-panels
\par\begingroup
\fontsize{10.5}{13.2}\selectfont
\begin{minipage}{\linewidth}\raggedright
\setlength{\LTleft}{0pt}\setlength{\LTright}{\fill}
\noindent
\begin{tabular}{@{}>{\raggedright\arraybackslash}p{22mm}>{\raggedright\arraybackslash}p{31.18mm}>{\raggedright\arraybackslash}p{31.18mm}>{\raggedright\arraybackslash}p{31.18mm}@{}}
\toprule
\textbf{Case} & \textbf{Masculine singular} & \textbf{Feminine singular} & \textbf{Neuter singular} \\
\midrule
Nominative & unrihtwīsa & unrihtwīse & unrihtwīse \\
Genitive & unrihtwīsan & unrihtwīsan & unrihtwīsan \\
Dative & unrihtwīsan & unrihtwīsan & unrihtwīsan \\
Accusative & unrihtwīsan & unrihtwīsan & unrihtwīse \\
Instrumental & unrihtwīsan & as dative & unrihtwīsan \\
\bottomrule
\end{tabular}
\par\addvspace{5pt}
\setlength{\LTleft}{0pt}\setlength{\LTright}{\fill}
\noindent
\begin{tabular}{@{}>{\raggedright\arraybackslash}p{22mm}>{\raggedright\arraybackslash}p{31.18mm}>{\raggedright\arraybackslash}p{31.18mm}>{\raggedright\arraybackslash}p{31.18mm}@{}}
\toprule
\textbf{Case} & \textbf{Masculine plural} & \textbf{Feminine plural} & \textbf{Neuter plural} \\
\midrule
Nominative & unrihtwīsan & unrihtwīsan & unrihtwīsan \\
Genitive & unrihtwīsra & unrihtwīsra & unrihtwīsra \\
Dative & unrihtwīsum & unrihtwīsum & unrihtwīsum \\
Accusative & unrihtwīsan & unrihtwīsan & unrihtwīsan \\
Instrumental & as dative & as dative & as dative \\
\bottomrule
\end{tabular}
\par\addvspace{5pt}
\end{minipage}\par\endgroup
% mdtable:T16:end

% mdparagraph:P050:begin
\textbf{Synfull}, “sinful”, likewise has no -u in those strong cells. Preserve the source headword's -ll in bare forms and before vowel endings. Before consonantal -ne, -re and -ra, simplify ll to l: \textbf{synfulne, synfulre, synfulra}. This spelling change affects the genitive plural in both declensions.
% mdparagraph:P050:end

% teaching: NEW-U10-ADJ-SYNFULL-STRONG
% table: NEW-U10-ADJ-SYNFULL-STRONG
% mdtable:T17:begin
% table-reflow: singular-plural-panels
\par\begingroup
\fontsize{10.5}{13.2}\selectfont
\begin{minipage}{\linewidth}\raggedright
\setlength{\LTleft}{0pt}\setlength{\LTright}{\fill}
\noindent
\begin{tabular}{@{}>{\raggedright\arraybackslash}p{22mm}>{\raggedright\arraybackslash}p{31.18mm}>{\raggedright\arraybackslash}p{31.18mm}>{\raggedright\arraybackslash}p{31.18mm}@{}}
\toprule
\textbf{Case} & \textbf{Masculine singular} & \textbf{Feminine singular} & \textbf{Neuter singular} \\
\midrule
Nominative & synfull & synfull & synfull \\
Genitive & synfulles & synfulre & synfulles \\
Dative & synfullum & synfulre & synfullum \\
Accusative & synfulne & synfulle & synfull \\
Instrumental & synfulle & as dative & synfulle \\
\bottomrule
\end{tabular}
\par\addvspace{5pt}
\setlength{\LTleft}{0pt}\setlength{\LTright}{\fill}
\noindent
\begin{tabular}{@{}>{\raggedright\arraybackslash}p{22mm}>{\raggedright\arraybackslash}p{31.18mm}>{\raggedright\arraybackslash}p{31.18mm}>{\raggedright\arraybackslash}p{31.18mm}@{}}
\toprule
\textbf{Case} & \textbf{Masculine plural} & \textbf{Feminine plural} & \textbf{Neuter plural} \\
\midrule
Nominative & synfulle & synfulla & synfull \\
Genitive & synfulra & synfulra & synfulra \\
Dative & synfullum & synfullum & synfullum \\
Accusative & synfulle & synfulla & synfull \\
Instrumental & as dative & as dative & as dative \\
\bottomrule
\end{tabular}
\par\addvspace{5pt}
\end{minipage}\par\endgroup
% mdtable:T17:end

% teaching: NEW-U10-ADJ-SYNFULL-WEAK
% table: NEW-U10-ADJ-SYNFULL-WEAK
% mdtable:T18:begin
% table-reflow: singular-plural-panels
\par\begingroup
\fontsize{10.5}{13.2}\selectfont
\begin{minipage}{\linewidth}\raggedright
\setlength{\LTleft}{0pt}\setlength{\LTright}{\fill}
\noindent
\begin{tabular}{@{}>{\raggedright\arraybackslash}p{22mm}>{\raggedright\arraybackslash}p{31.18mm}>{\raggedright\arraybackslash}p{31.18mm}>{\raggedright\arraybackslash}p{31.18mm}@{}}
\toprule
\textbf{Case} & \textbf{Masculine singular} & \textbf{Feminine singular} & \textbf{Neuter singular} \\
\midrule
Nominative & synfulla & synfulle & synfulle \\
Genitive & synfullan & synfullan & synfullan \\
Dative & synfullan & synfullan & synfullan \\
Accusative & synfullan & synfullan & synfulle \\
Instrumental & synfullan & as dative & synfullan \\
\bottomrule
\end{tabular}
\par\addvspace{5pt}
\setlength{\LTleft}{0pt}\setlength{\LTright}{\fill}
\noindent
\begin{tabular}{@{}>{\raggedright\arraybackslash}p{22mm}>{\raggedright\arraybackslash}p{31.18mm}>{\raggedright\arraybackslash}p{31.18mm}>{\raggedright\arraybackslash}p{31.18mm}@{}}
\toprule
\textbf{Case} & \textbf{Masculine plural} & \textbf{Feminine plural} & \textbf{Neuter plural} \\
\midrule
Nominative & synfullan & synfullan & synfullan \\
Genitive & synfulra & synfulra & synfulra \\
Dative & synfullum & synfullum & synfullum \\
Accusative & synfullan & synfullan & synfullan \\
Instrumental & as dative & as dative & as dative \\
\bottomrule
\end{tabular}
\par\addvspace{5pt}
\end{minipage}\par\endgroup
% mdtable:T18:end

\reserve{7\baselineskip}
% mdparagraph:P051:begin
\subsection*{The right side}
% mdparagraph:P051:end

% mdparagraph:P052:begin
\textbf{Swīðra}, “right” in the right-hand expression, has weak comparative morphology. Its stem is \textbf{swīðr-}. Use the complete weak table; no positive strong swīð paradigm is introduced here. In the genitive plural, the cluster from swīðr- plus -ra is simplified to \textbf{swīðra}, as in the earlier comparative pattern. The shape is a regular derivation under the cited cluster rule, not an additional strong adjective.
% mdparagraph:P052:end

% teaching: NEW-U10-ADJ-SWITHRA-WEAK
% table: NEW-U10-ADJ-SWITHRA-WEAK
% mdtable:T19:begin
% table-reflow: singular-plural-panels
\par\begingroup
\fontsize{10.5}{13.2}\selectfont
\begin{minipage}{\linewidth}\raggedright
\setlength{\LTleft}{0pt}\setlength{\LTright}{\fill}
\noindent
\begin{tabular}{@{}>{\raggedright\arraybackslash}p{22mm}>{\raggedright\arraybackslash}p{31.18mm}>{\raggedright\arraybackslash}p{31.18mm}>{\raggedright\arraybackslash}p{31.18mm}@{}}
\toprule
\textbf{Case} & \textbf{Masculine singular} & \textbf{Feminine singular} & \textbf{Neuter singular} \\
\midrule
Nominative & swīðra & swīðre & swīðre \\
Genitive & swīðran & swīðran & swīðran \\
Dative & swīðran & swīðran & swīðran \\
Accusative & swīðran & swīðran & swīðre \\
Instrumental & swīðran & as dative & swīðran \\
\bottomrule
\end{tabular}
\par\addvspace{5pt}
\setlength{\LTleft}{0pt}\setlength{\LTright}{\fill}
\noindent
\begin{tabular}{@{}>{\raggedright\arraybackslash}p{22mm}>{\raggedright\arraybackslash}p{31.18mm}>{\raggedright\arraybackslash}p{31.18mm}>{\raggedright\arraybackslash}p{31.18mm}@{}}
\toprule
\textbf{Case} & \textbf{Masculine plural} & \textbf{Feminine plural} & \textbf{Neuter plural} \\
\midrule
Nominative & swīðran & swīðran & swīðran \\
Genitive & swīðra & swīðra & swīðra \\
Dative & swīðrum & swīðrum & swīðrum \\
Accusative & swīðran & swīðran & swīðran \\
Instrumental & as dative & as dative & as dative \\
\bottomrule
\end{tabular}
\par\addvspace{5pt}
\end{minipage}\par\endgroup
% mdtable:T19:end

\reserve{7\baselineskip}
% mdparagraph:P053:begin
\subsection*{Adjective vocabulary and derivation}
% mdparagraph:P053:end

% mdparagraph:P054:begin
For \textbf{untrum, sellic, cwic} and \textbf{nacod}, replace the exact stem \textbf{hungrig-} by the listed stem in every cell of both hungrig tables. Keep the entire remainder of the model form. Retain the listed internal vowels; do not shorten untrum, sellic or nacod automatically. For \textbf{ēce}, replace \textbf{dēorwierð-} by \textbf{ēc-} in every corresponding cell of both dēorwierðe tables, retaining the final -e wherever the model has it. These operations include all genders, both numbers and both declensions; instrumental cells retain their reference limit.
% mdparagraph:P054:end

% teaching: NEW-U10-VOC-HUNGRIG
% diagnostic: NEW-U10-DIAG-HUNGRIG
% teaching: NEW-U10-VOC-DEORWIERTH
% diagnostic: NEW-U10-DIAG-DEORWIERTH
% teaching: NEW-U10-VOC-UNRIHTWIS
% diagnostic: NEW-U10-DIAG-UNRIHTWIS
% teaching: NEW-U10-VOC-SYNFULL
% diagnostic: NEW-U10-DIAG-SYNFULL
% teaching: NEW-U10-VOC-UNTRUM
% teaching: NEW-U10-ADJ-UNTRUM-STRONG
% teaching: NEW-U10-ADJ-UNTRUM-WEAK
% diagnostic: NEW-U10-DIAG-UNTRUM
% teaching: NEW-U10-VOC-SELLIC
% teaching: NEW-U10-ADJ-SELLIC-STRONG
% teaching: NEW-U10-ADJ-SELLIC-WEAK
% diagnostic: NEW-U10-DIAG-SELLIC
% teaching: NEW-U10-VOC-ECE
% teaching: NEW-U10-ADJ-ECE-STRONG
% teaching: NEW-U10-ADJ-ECE-WEAK
% diagnostic: NEW-U10-DIAG-ECE
% teaching: NEW-U10-VOC-CWIC
% teaching: NEW-U10-ADJ-CWIC-STRONG
% teaching: NEW-U10-ADJ-CWIC-WEAK
% diagnostic: NEW-U10-DIAG-CWIC
% teaching: NEW-U10-VOC-NACOD
% teaching: NEW-U10-ADJ-NACOD-STRONG
% teaching: NEW-U10-ADJ-NACOD-WEAK
% diagnostic: NEW-U10-DIAG-NACOD
% teaching: NEW-U10-VOC-SWITHRA
% diagnostic: NEW-U10-DIAG-SWITHRA
% mdtable:T20:begin
\par\begingroup
\fontsize{10.5}{13.2}\selectfont
\setlength{\LTleft}{0pt}\setlength{\LTright}{\fill}
\noindent
\begin{tabular}{@{}>{\raggedright\arraybackslash}p{46.15mm}>{\raggedright\arraybackslash}p{19mm}>{\raggedright\arraybackslash}p{53.22mm}@{}}
\toprule
\textbf{Headword and meaning} & \textbf{Stem} & \textbf{Complete model and operation} \\
\midrule
hungrig, hungry; hungry people & hungrig & Both full tables above; internal vowel retained. \\
dēorwierðe, precious; valuable & dēorwierð & Both full tables above; lexical -e replaced before other endings. \\
unrihtwīs, unrighteous; unrighteous people & unrihtwīs & Both full tables above; no -u in the long-stem cells. \\
synfull, sinful & synfull & Both full tables above; ll becomes l before -ne, -re, -ra. \\
untrum, sick; sick people & untrum & Replace hungrig- by untrum- in both tables. \\
sellic, strange; remarkable & sellic & Replace hungrig- by sellic- in both tables. \\
ēce, eternal & ēc & Replace dēorwierð- by ēc- in both tables. \\
cwic, alive; living & cwic & Replace hungrig- by cwic- in both tables. \\
nacod, naked; naked people & nacod & Replace hungrig- by nacod- in both tables. \\
swīðra, right, in a right-hand/\allowbreak{}side expression & swīðr & Full weak comparative table only; genitive plural swīðra. \\
\bottomrule
\end{tabular}
\par\endgroup

% mdtable:T20:end
% mdparagraph:P055:begin
For example, hungrigre gives \textbf{untrumre} at the same feminine singular genitive or dative coordinate; weak hungrigan gives \textbf{nacodan} at each corresponding coordinate. Dēorwierðe gives \textbf{ēce} where its lexical -e is retained, while dēorwierðum gives \textbf{ēcum}. Choosing a form still requires the case, gender, number and strong or weak use; matching one ending does not identify a unique cell.
% mdparagraph:P055:end

% teaching: NEW-U10-GRAM-u10-substantival-personal-adjectives
% mdparagraph:P056:begin
The article and a weak adjective can name people without an expressed noun: \textbf{the hungry, the sick, the naked, the unrighteous}. Use hungrig, untrum, nacod and unrihtwīs respectively. The whole phrase has the case required by its clause and the number intended. A plural group uses the plural adjective; a phrase describing one expressed person uses singular agreement.
% mdparagraph:P056:end

\reserve{8\baselineskip}
% example: NEW-U10-EXAMPLE14
% mdparagraph:P057:begin
Fēd ðā hungrigan.
% mdparagraph:P057:end

% mdparagraph:P058:begin
Feed the hungry.
% mdparagraph:P058:end

% mdparagraph:P059:begin
Fēd addresses one person. Ðā hungrigan is the accusative plural object, naming the hungry people. The number of the command belongs to its addressee, not to its object. With an expressed masculine singular noun, “the hungry stranger” instead has singular agreement, such as nominative \textbf{se hungriga cuma}.
% mdparagraph:P059:end

\reserve{7\baselineskip}
% mdparagraph:P060:begin
\section*{4. A member of a group and a quantity}
% mdparagraph:P060:end

% teaching: NEW-U10-GRAM-u10-partitive-group
% mdparagraph:P061:begin
A word naming a member or some members can take a \textbf{genitive plural whole}: ān, “one”, sum, “some”, and hwilc, “which”, can stand as heads followed by the group from which the selection is made. The head takes its case from the clause; the group remains genitive. Gender and number of the selecting word follow the intended member or members.
% mdparagraph:P061:end

% mdparagraph:P062:begin
Distinguish this from an agreeing word before an expressed noun. In \textbf{ān munuc}, “one monk”, ān modifies nominative singular munuc. In \textbf{ān ðāra munuca}, “one of the monks”, ān names the selected member and ðāra munuca is the genitive plural whole. This cardinal use of ān does not introduce the wider indefinite-article system.
% mdparagraph:P062:end

\reserve{8\baselineskip}
% example: NEW-U10-EXAMPLE04
% mdparagraph:P063:begin
Ān ðāra munuca cōm tō ðǣm hūse.
% mdparagraph:P063:end

% mdparagraph:P064:begin
One of the monks came to the house.
% mdparagraph:P064:end

% mdparagraph:P065:begin
Ān is nominative singular, the subject of cōm; ðāra munuca remains genitive plural. The same relation is available with plural sum, “some of”, and with interrogative hwilc, “which of”. A question can select one member or several; choose the corresponding singular or plural form and verb. These nominal genitive groups extend the earlier selected “of you/\allowbreak{}of them” phrases, without introducing every possible personal-pronoun genitive.
% mdparagraph:P065:end

% teaching: NEW-U10-GRAM-u10-fela-genitive
% teaching: NEW-U10-VOC-FELA
% mdparagraph:P066:begin
\textbf{Fela}, “many, a large number of”, is invariant in the selected construction and takes a genitive plural. Do not make it agree like an ordinary attributive adjective or a numeral.
% mdparagraph:P066:end

\reserve{8\baselineskip}
% example: NEW-U10-EXAMPLE03
% mdparagraph:P067:begin
Ic hæbbe fela scēapa.
% mdparagraph:P067:end

% mdparagraph:P068:begin
I have many sheep.
% mdparagraph:P068:end

% mdparagraph:P069:begin
Scēapa is genitive plural after fela. In the ordinary counted phrase \textbf{twā scēap}, “two sheep”, the numeral and noun instead share the case required by the clause. The genitive after fela is therefore not the form of scēap required directly by habban. Higher numerals and singular manig are outside the present construction.
% mdparagraph:P069:end

\reserve{7\baselineskip}
% mdparagraph:P070:begin
\section*{5. Datives and experiencers}
% mdparagraph:P070:end

% teaching: NEW-U10-GRAM-u10-dative-roles
% teaching: NEW-U10-GRAM-u10-don-beneficiary
% teaching: NEW-U10-VOC-DON
% mdparagraph:P071:begin
The dative has several functions. Identify the construction that requires it:
% mdparagraph:P071:end

\begin{itemize}
% mdparagraph:P072:begin
\item A \textbf{recipient} receives something: with sęllan, “give”, the thing is accusative and the recipient dative; tōgelǣdan, “bring”, follows the same case arrangement.
% mdparagraph:P072:end
% mdparagraph:P073:begin
\item A \textbf{beneficiary or affected person} is the person for or to whom a deed is done. In the selected use of dōn, “do”, the deed or “it” is accusative and the person dative.
% mdparagraph:P073:end
% mdparagraph:P074:begin
\item A \textbf{governed person} has the case required by a particular verb or adjective. Continue to use helpan, fręmman and behēfe within their learned meanings and dative frames. Forwiernan and unnan likewise require a dative person, although their thing is genitive.
% mdparagraph:P074:end
% mdparagraph:P075:begin
\item An \textbf{experiencer} is the person to whom something seems a certain way. The construction with ðyncan below uses a dative experiencer.
% mdparagraph:P075:end
\end{itemize}

\reserve{8\baselineskip}
% example: NEW-U10-EXAMPLE05
% mdparagraph:P076:begin
Se fiscere hilpð ðǣm cuman.
% mdparagraph:P076:end

% mdparagraph:P077:begin
The fisherman helps the stranger.
% mdparagraph:P077:end

% mdparagraph:P078:begin
The person helped is dative after hilpð. Its English translation has no preposition.
% mdparagraph:P078:end

\reserve{8\baselineskip}
% example: NEW-U10-EXAMPLE06
% mdparagraph:P079:begin
Sēo gāt is ðǣm scēaphierde behēfu.
% mdparagraph:P079:end

% mdparagraph:P080:begin
The goat is useful to the shepherd.
% mdparagraph:P080:end

% mdparagraph:P081:begin
The goat is the nominative subject. Behēfu agrees with feminine nominative singular gāt; ðǣm scēaphierde is dative, the person to whom the goat is useful. The adjective's agreement and the case of its complement are separate requirements.
% mdparagraph:P081:end

% teaching: NEW-U10-GRAM-u10-thyncan-experiencer
% teaching: NEW-U10-VOC-THYNCAN
% mdparagraph:P082:begin
In \textbf{ðyncan}, “seem”, the selected construction has a dative experiencer, a singular neuter nominative subject \textbf{hit, ðæt} or \textbf{ðis}, and the strong neuter predicate \textbf{gōd}. The finite verb is third-person singular: present ðyncð, preterite ðūhte. This is the complete present/\allowbreak{}preterite system needed for that defined construction.
% mdparagraph:P082:end

% teaching: NEW-U10-PAR-thyncan-EXPERIENCER-IND
% table: NEW-U10-PAR-thyncan-EXPERIENCER-IND
% mdtable:T21:begin
\par\begingroup
\fontsize{10.5}{13.2}\selectfont
\setlength{\LTleft}{0pt}\setlength{\LTright}{\fill}
\noindent
\begin{tabular}{@{}>{\raggedright\arraybackslash}p{42.37mm}>{\raggedright\arraybackslash}p{38.0mm}>{\raggedright\arraybackslash}p{38.0mm}@{}}
\toprule
\textbf{Person /\allowbreak{} number} & \textbf{Present} & \textbf{Preterite} \\
\midrule
3 singular & ðyncð & ðūhte \\
\bottomrule
\end{tabular}
\par\endgroup

% mdtable:T21:end
\reserve{8\baselineskip}
% example: NEW-U10-EXAMPLE07
% mdparagraph:P083:begin
Mē ðyncð ðis gōd.
% mdparagraph:P083:end

% mdparagraph:P084:begin
This seems good to me.
% mdparagraph:P084:end

% mdparagraph:P085:begin
Ðis is the nominative subject: “this” seems good. Mē is the dative experiencer. A plural experiencer does not make this singular subject plural. This use has an expressed nominative subject; it is not the subjectless construction that follows. Do not infer a general content-clause or passive construction from these forms.
% mdparagraph:P085:end

% teaching: NEW-U10-GRAM-u10-hyngran-impersonal
% teaching: NEW-U10-VOC-HYNGRAN
% mdparagraph:P086:begin
\textbf{Hyngran}, “be hungry”, is used here impersonally, with \textbf{no nominative subject}. The experiencer is accusative. The verb remains third-person singular even if the experiencer denotes several people: present hyngreð, preterite hyngrede. These two cells constitute the complete assigned impersonal indicative system; do not construct first- or second-person forms for it.
% mdparagraph:P086:end

% teaching: NEW-U10-PAR-hyngran-IMPERSONAL-IND
% table: NEW-U10-PAR-hyngran-IMPERSONAL-IND
% mdtable:T22:begin
\par\begingroup
\fontsize{10.5}{13.2}\selectfont
\setlength{\LTleft}{0pt}\setlength{\LTright}{\fill}
\noindent
\begin{tabular}{@{}>{\raggedright\arraybackslash}p{42.37mm}>{\raggedright\arraybackslash}p{38.0mm}>{\raggedright\arraybackslash}p{38.0mm}@{}}
\toprule
\textbf{Person /\allowbreak{} number} & \textbf{Present} & \textbf{Preterite} \\
\midrule
3 singular & hyngreð & hyngrede \\
\bottomrule
\end{tabular}
\par\endgroup

% mdtable:T22:end
\reserve{8\baselineskip}
% example: NEW-U10-EXAMPLE08
% mdparagraph:P087:begin
Hine hyngrede.
% mdparagraph:P087:end

% mdparagraph:P088:begin
He was hungry.
% mdparagraph:P088:end

% mdparagraph:P089:begin
Hine is accusative, although English requires “he” as its subject. Compare the ordinary copular construction \textbf{Hē wæs hungrig}, “he was hungry”: hē is nominative, wæs agrees with it, and hungrig is a predicate adjective. Both express hunger, but their grammar differs. The accusative experiencer pattern is the one to use when a task expressly requires impersonal hyngran.
% mdparagraph:P089:end

\reserve{7\baselineskip}
% mdparagraph:P090:begin
\section*{6. Accompaniment, means and the instrumental}
% mdparagraph:P090:end

% teaching: NEW-U10-VOC-MID
% teaching: NEW-U10-GRAM-u10-mid-means
% teaching: NEW-U10-GRAM-u10-fedan-accusative-means
% mdparagraph:P091:begin
\textbf{Mid + dative} can express accompaniment, “with”, or means, “with, by means of”. The preposition governs the dative in both uses. The relation must be read from the situation.
% mdparagraph:P091:end

% mdparagraph:P092:begin
\textbf{Fēdan}, “feed, nourish”, takes the person or animal fed in the accusative. Food used for feeding can be expressed by mid + dative.
% mdparagraph:P092:end

\reserve{8\baselineskip}
% example: NEW-U10-EXAMPLE09
% mdparagraph:P093:begin
Se munuc cōm mid ðǣm cuman.
% mdparagraph:P093:end

% mdparagraph:P094:begin
The monk came with the stranger.
% mdparagraph:P094:end

% mdparagraph:P095:begin
Here mid names the accompanying person.
% mdparagraph:P095:end

\reserve{8\baselineskip}
% example: NEW-U10-EXAMPLE10
% mdparagraph:P096:begin
Se munuc fēdde ðone cuman mid gōdum męte.
% mdparagraph:P096:end

% mdparagraph:P097:begin
The monk fed the stranger with good food.
% mdparagraph:P097:end

% mdparagraph:P098:begin
Here mid names the means of feeding. Ðone cuman is the accusative person fed; gōdum męte is dative after mid. This is the ordinary prose means construction.
% mdparagraph:P098:end

% teaching: NEW-U10-GRAM-u10-bare-instrument
% teaching: NEW-U10-GRAM-u10-instrumental-morphology
% mdparagraph:P099:begin
Means may also be expressed by a \textbf{bare instrumental}, without a preposition. This use is especially characteristic of poetry; mid + dative is the usual prose choice. The present exercises deliberately practise the bare construction in the feeding expression.
% mdparagraph:P099:end

% mdparagraph:P100:begin
The demonstrative has distinct masculine/\allowbreak{}neuter singular instrumental \textbf{ðȳ}, with variant \textbf{ðon}. Strong gōd has instrumental \textbf{gōde} in those two genders. Compare their dative forms ðǣm and gōdum. Only these distinct instrumental forms are newly required for composition here; the previously displayed instruments of ðes and hwā remain reference.
% mdparagraph:P100:end

% teaching: NEW-U10-PAR-se-SELECTED-INSTR
% table: NEW-U10-PAR-se-SELECTED-INSTR
% mdtable:T23:begin
\par\begingroup
\fontsize{10.5}{13.2}\selectfont
\setlength{\LTleft}{0pt}\setlength{\LTright}{\fill}
\noindent
\begin{tabular}{@{}>{\raggedright\arraybackslash}p{38.0mm}>{\raggedright\arraybackslash}p{34.94mm}>{\raggedright\arraybackslash}p{45.43mm}@{}}
\toprule
\textbf{Gender /\allowbreak{} number} & \textbf{Instrumental} & \textbf{Instrumental variant} \\
\midrule
Masculine singular & ðȳ & ðon \\
Neuter singular & ðȳ & ðon \\
\bottomrule
\end{tabular}
\par\endgroup

% mdtable:T23:end
% teaching: NEW-U10-PAR-god-adjective-SELECTED-INSTR
% table: NEW-U10-PAR-god-adjective-SELECTED-INSTR
% mdtable:T24:begin
\par\begingroup
\fontsize{10.5}{13.2}\selectfont
\setlength{\LTleft}{0pt}\setlength{\LTright}{\fill}
\noindent
\begin{tabular}{@{}>{\raggedright\arraybackslash}p{56.08mm}>{\raggedright\arraybackslash}p{65.11mm}@{}}
\toprule
\textbf{Gender /\allowbreak{} number} & \textbf{Strong instrumental} \\
\midrule
Masculine singular & gōde \\
Neuter singular & gōde \\
\bottomrule
\end{tabular}
\par\endgroup

% mdtable:T24:end
% mdparagraph:P101:begin
Nouns use the dative-shaped form for this instrumental relation: \textbf{męte} does not acquire a new ending. Feminine singular and all plural modifiers likewise share their forms with the dative. Thus the feminine singular forms are \textbf{ðǣre} and strong \textbf{gōdre}, and the plural forms \textbf{ðǣm} and \textbf{gōdum}. Shared shape does not remove the distinction of function. In a weak adjective, use the corresponding weak dative shape; the new adjectives' separate instrumental cells remain reference.
% mdparagraph:P101:end

\reserve{8\baselineskip}
% example: NEW-U10-EXAMPLE11
% mdparagraph:P102:begin
Se munuc fēdde ðone cuman gōde męte.
% mdparagraph:P102:end

% mdparagraph:P103:begin
The monk fed the stranger with good food.
% mdparagraph:P103:end

% mdparagraph:P104:begin
Gōde is a distinct strong masculine singular instrumental. Męte has the same shape as the dative singular, but the phrase expresses bare means.
% mdparagraph:P104:end

\reserve{8\baselineskip}
% example: NEW-U10-EXAMPLE12
% mdparagraph:P105:begin
Se munuc fēdde ðone cuman ðȳ męte.
% mdparagraph:P105:end

% mdparagraph:P106:begin
The monk fed the stranger with that food.
% mdparagraph:P106:end

% mdparagraph:P107:begin
Ðȳ is the demonstrative instrumental agreeing with masculine męte. Adding mid changes the construction: use the dative \textbf{mid ðǣm męte}. Do not treat mid + a distinct instrumental as the ordinary prose pattern taught here. An instrumental meaning, a distinct morphological form, and the case governed by a preposition are three separate facts.
% mdparagraph:P107:end

\reserve{7\baselineskip}
% mdparagraph:P108:begin
\section*{7. Complete verb systems and their constructions}
% mdparagraph:P108:end

% mdparagraph:P109:begin
The following tables give complete present and preterite indicative systems, except where the expressly defined impersonal or experiencer system has already been given. Commands are included only for the verbs with a command table. Unnan has its supported indicative and complete subjunctive system. Infinitives used as dictionary headwords identify verbs; this alone does not assign every infinitive construction. The newly usable plain infinitives are specified in §8.
% mdparagraph:P109:end

% mdparagraph:P110:begin
Plural subjects use the common plural ending in each indicative tense; dual subjects also use the plural finite form. A direct command's number belongs to the addressee. The object's number does not determine it.
% mdparagraph:P110:end

\reserve{7\baselineskip}
% mdparagraph:P111:begin
\subsection*{Weak class I: stem changes and dental preterites}
% mdparagraph:P111:end

% mdparagraph:P112:begin
The ordinary selected present endings are singular -e, -st, -ð and plural -að, with the contractions printed below. The preterite uses a dental stem followed by singular -e, -est, -e and plural -on. A long vowel followed by d gives the dental sequence \textbf{dd} in fēdde, scrȳdde and tōgelǣdde. Their second- and third-singular present forms instead have \textbf{tst} and \textbf{t}. Keep the long vowel in these selected forms.
% mdparagraph:P112:end

% teaching: NEW-U10-VOC-FEDAN
% teaching: NEW-U10-GRAM-u10-fedan-indicative
% teaching: NEW-U10-GRAM-u10-fedan-commands
% mdparagraph:P113:begin
\textbf{Fēdan}, “feed, nourish”, has present stem fēd- and preterite stem fēdd-. It feeds an accusative person or animal; its means expressions were given in §6. The singular command is the bare stem fēd; the plural command equals fēdað.
% mdparagraph:P113:end

\reserve{7\baselineskip}
% mdparagraph:P114:begin
\subsection*{Fēdan: indicative}
% mdparagraph:P114:end

% teaching: NEW-U10-PAR-fedan-IND
% table: NEW-U10-PAR-fedan-IND
% mdtable:T25:begin
\par\begingroup
\fontsize{10.5}{13.2}\selectfont
\setlength{\LTleft}{0pt}\setlength{\LTright}{\fill}
\noindent
\begin{tabular}{@{}>{\raggedright\arraybackslash}p{42.37mm}>{\raggedright\arraybackslash}p{38.0mm}>{\raggedright\arraybackslash}p{38.0mm}@{}}
\toprule
\textbf{Person /\allowbreak{} number} & \textbf{Present} & \textbf{Preterite} \\
\midrule
1 singular & fēde & fēdde \\
2 singular & fētst & fēddest \\
3 singular & fēt & fēdde \\
1 plural & fēdað & fēddon \\
2 plural & fēdað & fēddon \\
3 plural & fēdað & fēddon \\
\bottomrule
\end{tabular}
\par\endgroup

% mdtable:T25:end
\reserve{7\baselineskip}
% mdparagraph:P115:begin
\subsection*{Fēdan: commands}
% mdparagraph:P115:end

% teaching: NEW-U10-PAR-fedan-IMP
% table: NEW-U10-PAR-fedan-IMP
% mdtable:T26:begin
\par\begingroup
\fontsize{10.5}{13.2}\selectfont
\setlength{\LTleft}{0pt}\setlength{\LTright}{\fill}
\noindent
\begin{tabular}{@{}>{\raggedright\arraybackslash}p{60.59mm}>{\raggedright\arraybackslash}p{60.59mm}@{}}
\toprule
\textbf{Addressee} & \textbf{Command} \\
\midrule
Singular & fēd \\
Plural & fēdað \\
\bottomrule
\end{tabular}
\par\endgroup

% mdtable:T26:end
% teaching: NEW-U10-VOC-SCRYDAN
% teaching: NEW-U10-GRAM-u10-scrydan-indicative
% teaching: NEW-U10-GRAM-u10-scrydan-commands
% teaching: NEW-U10-GRAM-u10-scrydan-accusative
% mdparagraph:P116:begin
\textbf{Scrȳdan}, “clothe, dress”, takes an accusative person. Its present stem is scrȳd-, its preterite stem scrȳdd-. The commands are scrȳd and scrȳdað. The assigned construction names the person clothed; a further garment complement is not required in this unit.
% mdparagraph:P116:end

\reserve{7\baselineskip}
% mdparagraph:P117:begin
\subsection*{Scrȳdan: indicative}
% mdparagraph:P117:end

% teaching: NEW-U10-PAR-scrydan-IND
% table: NEW-U10-PAR-scrydan-IND
% mdtable:T27:begin
\par\begingroup
\fontsize{10.5}{13.2}\selectfont
\setlength{\LTleft}{0pt}\setlength{\LTright}{\fill}
\noindent
\begin{tabular}{@{}>{\raggedright\arraybackslash}p{42.37mm}>{\raggedright\arraybackslash}p{38.0mm}>{\raggedright\arraybackslash}p{38.0mm}@{}}
\toprule
\textbf{Person /\allowbreak{} number} & \textbf{Present} & \textbf{Preterite} \\
\midrule
1 singular & scrȳde & scrȳdde \\
2 singular & scrȳtst & scrȳddest \\
3 singular & scrȳt & scrȳdde \\
1 plural & scrȳdað & scrȳddon \\
2 plural & scrȳdað & scrȳddon \\
3 plural & scrȳdað & scrȳddon \\
\bottomrule
\end{tabular}
\par\endgroup

% mdtable:T27:end
\reserve{7\baselineskip}
% mdparagraph:P118:begin
\subsection*{Scrȳdan: commands}
% mdparagraph:P118:end

% teaching: NEW-U10-PAR-scrydan-IMP
% table: NEW-U10-PAR-scrydan-IMP
% mdtable:T28:begin
\par\begingroup
\fontsize{10.5}{13.2}\selectfont
\setlength{\LTleft}{0pt}\setlength{\LTright}{\fill}
\noindent
\begin{tabular}{@{}>{\raggedright\arraybackslash}p{60.59mm}>{\raggedright\arraybackslash}p{60.59mm}@{}}
\toprule
\textbf{Addressee} & \textbf{Command} \\
\midrule
Singular & scrȳd \\
Plural & scrȳdað \\
\bottomrule
\end{tabular}
\par\endgroup

% mdtable:T28:end
% teaching: NEW-U10-VOC-TOGELAEDAN
% teaching: NEW-U10-GRAM-u10-togelaedan-indicative
% teaching: NEW-U10-GRAM-u10-togelaedan-recipient
% mdparagraph:P119:begin
\textbf{Tōgelǣdan}, “bring”, takes an accusative thing and a dative recipient. Its present stem tōgelǣd- becomes tōgelǣdd- in the preterite. Use this verb when bringing something to someone in the assigned exercises. Its complete indicative is now available; no new command forms are assigned.
% mdparagraph:P119:end

\reserve{7\baselineskip}
% mdparagraph:P120:begin
\subsection*{Tōgelǣdan: indicative}
% mdparagraph:P120:end

% teaching: NEW-U10-PAR-togelaedan-IND
% table: NEW-U10-PAR-togelaedan-IND
% mdtable:T29:begin
\par\begingroup
\fontsize{10.5}{13.2}\selectfont
\setlength{\LTleft}{0pt}\setlength{\LTright}{\fill}
\noindent
\begin{tabular}{@{}>{\raggedright\arraybackslash}p{41.05mm}>{\raggedright\arraybackslash}p{36.82mm}>{\raggedright\arraybackslash}p{40.5mm}@{}}
\toprule
\textbf{Person /\allowbreak{} number} & \textbf{Present} & \textbf{Preterite} \\
\midrule
1 singular & tōgelǣde & tōgelǣdde \\
2 singular & tōgelǣtst & tōgelǣddest \\
3 singular & tōgelǣt & tōgelǣdde \\
1 plural & tōgelǣdað & tōgelǣddon \\
2 plural & tōgelǣdað & tōgelǣddon \\
3 plural & tōgelǣdað & tōgelǣddon \\
\bottomrule
\end{tabular}
\par\endgroup

% mdtable:T29:end
% teaching: NEW-U10-GRAM-u10-forwiernan-indicative
% mdparagraph:P121:begin
\textbf{Forwiernan}, “deny, refuse”, has present stem forwiern- and dental preterite stem forwiernd-. Preserve its dative person and genitive thing in every tense and person. The present third singular is forwiernð; the preterite third singular is forwiernde.
% mdparagraph:P121:end

\reserve{7\baselineskip}
% mdparagraph:P122:begin
\subsection*{Forwiernan: indicative}
% mdparagraph:P122:end

% teaching: NEW-U10-PAR-forwiernan-IND
% table: NEW-U10-PAR-forwiernan-IND
% mdtable:T30:begin
\par\begingroup
\fontsize{10.5}{13.2}\selectfont
\setlength{\LTleft}{0pt}\setlength{\LTright}{\fill}
\noindent
\begin{tabular}{@{}>{\raggedright\arraybackslash}p{39.82mm}>{\raggedright\arraybackslash}p{35.71mm}>{\raggedright\arraybackslash}p{42.85mm}@{}}
\toprule
\textbf{Person /\allowbreak{} number} & \textbf{Present} & \textbf{Preterite} \\
\midrule
1 singular & forwierne & forwiernde \\
2 singular & forwiernst & forwierndest \\
3 singular & forwiernð & forwiernde \\
1 plural & forwiernað & forwierndon \\
2 plural & forwiernað & forwierndon \\
3 plural & forwiernað & forwierndon \\
\bottomrule
\end{tabular}
\par\endgroup

% mdtable:T30:end
% teaching: NEW-U10-VOC-SCENCAN
% teaching: NEW-U10-GRAM-u10-scencan-indicative
% teaching: NEW-U10-GRAM-u10-scencan-recipient
% mdparagraph:P123:begin
\textbf{Scęncan}, “pour out, give to drink”, takes an accusative drink and a dative recipient. It retains ę and the present stem scęnc-. Its voiceless consonant is followed by \textbf{t} in the preterite stem scęnct-, giving scęncte, not a form in -de.
% mdparagraph:P123:end

\reserve{7\baselineskip}
% mdparagraph:P124:begin
\subsection*{Scęncan: indicative}
% mdparagraph:P124:end

% teaching: NEW-U10-PAR-scencan-IND
% table: NEW-U10-PAR-scencan-IND
% mdtable:T31:begin
\par\begingroup
\fontsize{10.5}{13.2}\selectfont
\setlength{\LTleft}{0pt}\setlength{\LTright}{\fill}
\noindent
\begin{tabular}{@{}>{\raggedright\arraybackslash}p{42.37mm}>{\raggedright\arraybackslash}p{38.0mm}>{\raggedright\arraybackslash}p{38.0mm}@{}}
\toprule
\textbf{Person /\allowbreak{} number} & \textbf{Present} & \textbf{Preterite} \\
\midrule
1 singular & scęnce & scęncte \\
2 singular & scęncst & scęnctest \\
3 singular & scęncð & scęncte \\
1 plural & scęncað & scęncton \\
2 plural & scęncað & scęncton \\
3 plural & scęncað & scęncton \\
\bottomrule
\end{tabular}
\par\endgroup

% mdtable:T31:end
% teaching: NEW-U10-VOC-WYRCEAN
% teaching: NEW-U10-GRAM-u10-wyrcean-work
% teaching: NEW-U10-GRAM-u10-wyrcean-commands
% mdparagraph:P125:begin
The familiar \textbf{wyrcean} now also means “work, perform”, with an accusative work or deed, including cognate \textbf{weorc}, “work”. This ordinary sense differs from its earlier yield-fruit use. Its indicative forms are already learned. The two direct commands are given here; retain the different stem spellings wyrc- and wyrce-.
% mdparagraph:P125:end

\reserve{7\baselineskip}
% mdparagraph:P126:begin
\subsection*{Wyrcean: commands}
% mdparagraph:P126:end

% teaching: NEW-U10-PAR-wyrcean-IMP
% table: NEW-U10-PAR-wyrcean-IMP
% mdtable:T32:begin
\par\begingroup
\fontsize{10.5}{13.2}\selectfont
\setlength{\LTleft}{0pt}\setlength{\LTright}{\fill}
\noindent
\begin{tabular}{@{}>{\raggedright\arraybackslash}p{60.59mm}>{\raggedright\arraybackslash}p{60.59mm}@{}}
\toprule
\textbf{Addressee} & \textbf{Command} \\
\midrule
Singular & wyrc \\
Plural & wyrceað \\
\bottomrule
\end{tabular}
\par\endgroup

% mdtable:T32:end
\reserve{7\baselineskip}
% mdparagraph:P127:begin
\subsection*{Weak class II}
% mdparagraph:P127:end

% mdparagraph:P128:begin
The selected class II endings are present -ie, -ast, -að, plural -iað; preterite -ode, -odest, -ode, plural -odon. The class identifies the inflection, not the case of a verb's complement. Cook also records spellings with -ig(e)- or -g- and preterite -edon; these remain reference variants, not extra forms to reproduce here.
% mdparagraph:P128:end

% teaching: NEW-U10-VOC-ANDSWARIAN
% teaching: NEW-U10-GRAM-u10-andswarian-indicative
% teaching: NEW-U10-GRAM-u10-andswarian-dative
% mdparagraph:P129:begin
\textbf{Andswarian}, “answer, reply”, takes a \textbf{dative addressee}. A direct quotation can give the reply. Use the weak class II stem andswar- in this table; the separate weak class I variant andswerian is not introduced.
% mdparagraph:P129:end

\reserve{7\baselineskip}
% mdparagraph:P130:begin
\subsection*{Andswarian: indicative}
% mdparagraph:P130:end

% teaching: NEW-U10-PAR-andswarian-IND
% table: NEW-U10-PAR-andswarian-IND
% mdtable:T33:begin
\par\begingroup
\fontsize{10.5}{13.2}\selectfont
\setlength{\LTleft}{0pt}\setlength{\LTright}{\fill}
\noindent
\begin{tabular}{@{}>{\raggedright\arraybackslash}p{39.82mm}>{\raggedright\arraybackslash}p{35.71mm}>{\raggedright\arraybackslash}p{42.85mm}@{}}
\toprule
\textbf{Person /\allowbreak{} number} & \textbf{Present} & \textbf{Preterite} \\
\midrule
1 singular & andswarie & andswarode \\
2 singular & andswarast & andswarodest \\
3 singular & andswarað & andswarode \\
1 plural & andswariað & andswarodon \\
2 plural & andswariað & andswarodon \\
3 plural & andswariað & andswarodon \\
\bottomrule
\end{tabular}
\par\endgroup

% mdtable:T33:end
% teaching: NEW-U10-VOC-GENEOSIAN
% teaching: NEW-U10-GRAM-u10-geneosian-indicative
% teaching: NEW-U10-GRAM-u10-geneosian-commands
% teaching: NEW-U10-GRAM-u10-geneosian-accusative
% mdparagraph:P131:begin
\textbf{Genēosian}, “visit”, takes an \textbf{accusative person} in the selected construction. Keep Cook's long ē in genēos-. The singular command is genēosa and the plural genēosiað. Do not transfer the genitive government of the related unprefixed nēosian to this selected prefixed use.
% mdparagraph:P131:end

\reserve{7\baselineskip}
% mdparagraph:P132:begin
\subsection*{Genēosian: indicative}
% mdparagraph:P132:end

% teaching: NEW-U10-PAR-geneosian-IND
% table: NEW-U10-PAR-geneosian-IND
% mdtable:T34:begin
\par\begingroup
\fontsize{10.5}{13.2}\selectfont
\setlength{\LTleft}{0pt}\setlength{\LTright}{\fill}
\noindent
\begin{tabular}{@{}>{\raggedright\arraybackslash}p{41.05mm}>{\raggedright\arraybackslash}p{36.82mm}>{\raggedright\arraybackslash}p{40.5mm}@{}}
\toprule
\textbf{Person /\allowbreak{} number} & \textbf{Present} & \textbf{Preterite} \\
\midrule
1 singular & genēosie & genēosode \\
2 singular & genēosast & genēosodest \\
3 singular & genēosað & genēosode \\
1 plural & genēosiað & genēosodon \\
2 plural & genēosiað & genēosodon \\
3 plural & genēosiað & genēosodon \\
\bottomrule
\end{tabular}
\par\endgroup

% mdtable:T34:end
\reserve{7\baselineskip}
% mdparagraph:P133:begin
\subsection*{Genēosian: commands}
% mdparagraph:P133:end

% teaching: NEW-U10-PAR-geneosian-IMP
% table: NEW-U10-PAR-geneosian-IMP
% mdtable:T35:begin
\par\begingroup
\fontsize{10.5}{13.2}\selectfont
\setlength{\LTleft}{0pt}\setlength{\LTright}{\fill}
\noindent
\begin{tabular}{@{}>{\raggedright\arraybackslash}p{60.59mm}>{\raggedright\arraybackslash}p{60.59mm}@{}}
\toprule
\textbf{Addressee} & \textbf{Command} \\
\midrule
Singular & genēosa \\
Plural & genēosiað \\
\bottomrule
\end{tabular}
\par\endgroup

% mdtable:T35:end
\reserve{7\baselineskip}
% mdparagraph:P134:begin
\subsection*{Strong class I: gewītan}
% mdparagraph:P134:end

% teaching: NEW-U10-VOC-GEWITAN
% teaching: NEW-U10-GRAM-u10-gewitan-indicative
% teaching: NEW-U10-GRAM-u10-gewitan-departure
% mdparagraph:P135:begin
\textbf{Gewītan}, “depart, go away”, can be followed by \textbf{fram + dative}, naming the person or place left. This strong class I verb has long ī in the present, ā in the first and third singular preterite, and short i in the second singular and the plural preterite. The t-stem contracts in gewītst and gewīt; Cook's tt/\allowbreak{}t variation is a reference spelling variation. The forms below select single t.
% mdparagraph:P135:end

% mdparagraph:P136:begin
Distinguish this departure verb from the different verb meaning “know”. An accompanying reflexive dative is attested with departure, but is not required in the selected construction.
% mdparagraph:P136:end

\reserve{7\baselineskip}
% mdparagraph:P137:begin
\subsection*{Gewītan: indicative}
% mdparagraph:P137:end

% teaching: NEW-U10-PAR-gewitan-IND
% table: NEW-U10-PAR-gewitan-IND
% mdtable:T36:begin
\par\begingroup
\fontsize{10.5}{13.2}\selectfont
\setlength{\LTleft}{0pt}\setlength{\LTright}{\fill}
\noindent
\begin{tabular}{@{}>{\raggedright\arraybackslash}p{42.37mm}>{\raggedright\arraybackslash}p{38.0mm}>{\raggedright\arraybackslash}p{38.0mm}@{}}
\toprule
\textbf{Person /\allowbreak{} number} & \textbf{Present} & \textbf{Preterite} \\
\midrule
1 singular & gewīte & gewāt \\
2 singular & gewītst & gewite \\
3 singular & gewīt & gewāt \\
1 plural & gewītað & gewiton \\
2 plural & gewītað & gewiton \\
3 plural & gewītað & gewiton \\
\bottomrule
\end{tabular}
\par\endgroup

% mdtable:T36:end
\reserve{7\baselineskip}
% mdparagraph:P138:begin
\subsection*{Unnan: supported forms, wishes and commands}
% mdparagraph:P138:end

% teaching: NEW-U10-GRAM-u10-unnan-indicative
% mdparagraph:P139:begin
\textbf{Unnan}, “grant, allow”, is a preterite-present verb: the present indicative has ann in the first and third singular and unnon throughout the plural. Its dental preterite has ūð-. Remember its dative person and genitive thing.
% mdparagraph:P139:end

% mdparagraph:P140:begin
The checked sources do not supply a second-singular present indicative form for this table. The dash is an \textbf{evidence gap}, not a form and not a claim that the form could never occur. Do not fill it by analogy.
% mdparagraph:P140:end

\reserve{7\baselineskip}
% mdparagraph:P141:begin
\subsection*{Unnan: indicative}
% mdparagraph:P141:end

% teaching: NEW-U10-PAR-unnan-IND
% table: NEW-U10-PAR-unnan-IND
% mdtable:T37:begin
\par\begingroup
\fontsize{10.5}{13.2}\selectfont
\setlength{\LTleft}{0pt}\setlength{\LTright}{\fill}
\noindent
\begin{tabular}{@{}>{\raggedright\arraybackslash}p{42.37mm}>{\raggedright\arraybackslash}p{38.0mm}>{\raggedright\arraybackslash}p{38.0mm}@{}}
\toprule
\textbf{Person /\allowbreak{} number} & \textbf{Present} & \textbf{Preterite} \\
\midrule
1 singular & ann & ūðe \\
2 singular & — & ūðest \\
3 singular & ann & ūðe \\
1 plural & unnon & ūðon \\
2 plural & unnon & ūðon \\
3 plural & unnon & ūðon \\
\bottomrule
\end{tabular}
\par\endgroup

% mdtable:T37:end
% teaching: NEW-U10-GRAM-u10-unnan-subjunctive-command
% mdparagraph:P142:begin
The present subjunctive is \textbf{unne} in every singular person and \textbf{unnen} in every plural person. The preterite subjunctive is \textbf{ūðe} throughout the singular and \textbf{ūðen} throughout the plural. These forms serve the earlier uses of subjunctive wishes: a present grant or wish uses unne/\allowbreak{}unnen; an unreal present wish can use the preterite subjunctive. Case government remains unchanged.
% mdparagraph:P142:end

\reserve{7\baselineskip}
% mdparagraph:P143:begin
\subsection*{Unnan: subjunctive}
% mdparagraph:P143:end

% teaching: NEW-U10-PAR-unnan-SUBJ
% table: NEW-U10-PAR-unnan-SUBJ
% mdtable:T38:begin
\par\begingroup
\fontsize{10.5}{13.2}\selectfont
\setlength{\LTleft}{0pt}\setlength{\LTright}{\fill}
\noindent
\begin{tabular}{@{}>{\raggedright\arraybackslash}p{42.37mm}>{\raggedright\arraybackslash}p{38.0mm}>{\raggedright\arraybackslash}p{38.0mm}@{}}
\toprule
\textbf{Person /\allowbreak{} number} & \textbf{Present} & \textbf{Preterite} \\
\midrule
1 singular & unne & ūðe \\
2 singular & unne & ūðe \\
3 singular & unne & ūðe \\
1 plural & unnen & ūðen \\
2 plural & unnen & ūðen \\
3 plural & unnen & ūðen \\
\bottomrule
\end{tabular}
\par\endgroup

% mdtable:T38:end
% mdparagraph:P144:begin
For a direct addressed command “grant”, use \textbf{unne} to one addressee and \textbf{unnen} to several. These are the second-person subjunctive forms used as command substitutes. Cook labels unne “Imper.”; Fulk describes the subjunctive substitution. Both accounts support the same singular surface. Do not make a formal plural command in -að, or mistake command unne for the missing indicative cell. The preterite plural indicative ūðon and subjunctive ūðen differ in their final vowel.
% mdparagraph:P144:end

\reserve{7\baselineskip}
% mdparagraph:P145:begin
\section*{8. Further constructions for the exercises}
% mdparagraph:P145:end

% teaching: NEW-U10-GRAM-u10-perception-object-predicate
% mdparagraph:P146:begin
\textbf{Gesēon}, “see”, can take an accusative object followed by an adjective describing the object's perceived condition. The predicate adjective is strong and agrees with the object in case, gender and number. It does not describe the observer.
% mdparagraph:P146:end

\reserve{8\baselineskip}
% example: NEW-U10-EXAMPLE13
% mdparagraph:P147:begin
Ic geseah ðone cuman untrumne.
% mdparagraph:P147:end

% mdparagraph:P148:begin
I saw the stranger sick.
% mdparagraph:P148:end

% mdparagraph:P149:begin
Untrumne is masculine accusative singular, agreeing with ðone cuman. It is a description of the object, distinct from the earlier accusative actor followed by an infinitive naming an action.
% mdparagraph:P149:end

% teaching: NEW-U10-GRAM-u10-hieran-perception-infinitive
% teaching: NEW-U10-GRAM-u10-promoted-plain-infinitives
% teaching: NEW-U10-VOC-HIERAN
% teaching: NEW-U10-VOC-BLAWAN
% teaching: NEW-U10-VOC-CIEPAN
% teaching: NEW-U10-VOC-BYCGAN
% mdparagraph:P150:begin
\textbf{Hīeran}, “hear”, now takes the same defined action construction as gesēon: an \textbf{accusative perceived actor + plain infinitive}. The infinitive names what the actor is heard doing and has no preceding tō. \textbf{Blāwan}, “blow”, is available as that plain infinitive. If the actor is a noun denoting an instrument or implement, it is still accusative in this construction.
% mdparagraph:P150:end

% mdparagraph:P151:begin
The plain infinitives \textbf{cīepan}, “sell”, and \textbf{bycgan}, “buy”, are also available after volitional \textbf{willan}. The person who wills is the understood subject of the infinitive. Keep each infinitive's ordinary accusative thing sold or bought. This use concerns willingness or intention; it does not establish a mechanical future tense.
% mdparagraph:P151:end

% teaching: NEW-U10-GRAM-u10-present-future
% mdparagraph:P152:begin
An Old English present can refer to a future event when the context makes the time clear. English “will come”, “will see” or “will separate” can therefore be translated with the ordinary present. \textbf{Willan + infinitive} is appropriate when intention or willingness is actually meant, including the farmer's and workman's stated plans. English “will” alone does not require willan; likewise “shall come” can predict an event without commanding it.
% mdparagraph:P152:end

% mdparagraph:P153:begin
Keep the earlier negative volitional use: ne with willan expresses unwillingness or refusal. Use that meaning when it suits the speaker's refusal to act. Where a prompt admits either prediction or intention, the construction must express the chosen meaning consistently.
% mdparagraph:P153:end

% teaching: NEW-U10-GRAM-u10-to-destination
% teaching: NEW-U10-GRAM-u10-beforan-location
% teaching: NEW-U10-VOC-TO
% teaching: NEW-U10-VOC-BEFORAN
% mdparagraph:P154:begin
\textbf{Tō + dative} names a destination with verbs of coming, going or flying. It is distinct from a purpose/\allowbreak{}function expression and from tō before an inflected infinitive. \textbf{Beforan + dative} means “before, in front of” a person in the standing construction. The exact reading below uses the separately supplied ætforan; do not silently replace its spelling.
% mdparagraph:P154:end

% teaching: NEW-U10-GRAM-u10-on-side-idiom
% teaching: NEW-U10-GRAM-u10-on-name-idiom
% teaching: NEW-U10-VOC-ON
% mdparagraph:P155:begin
In the bounded expression \textbf{on … swīðran hand}, “on … right hand/\allowbreak{}side”, the noun \textbf{hand is accusative}. The expression occurs even with static standan, “stand”. Learn this case with the side-expression; motion alone does not explain it. The bare noun hand distinguishes the selected accusative from dative hande or handa. Weak swīðran by itself cannot decide the case, since -an occurs in several weak singular cells. The parallel left-side expression in the reading is supplied locally.
% mdparagraph:P155:end

% mdparagraph:P156:begin
Ordinary \textbf{on + dative} still expresses location, as with a seat, prison or inn. The side-expression does not replace that general pattern.
% mdparagraph:P156:end

% mdparagraph:P157:begin
\textbf{On … naman}, “in … name”, instead has dative \textbf{naman}. A possessor between on and naman is genitive. As with the recipient phrase in §1, distinguish the case governed by the preposition from the genitive inside its complement.
% mdparagraph:P157:end

% teaching: NEW-U10-GRAM-u10-totwaeman-separation
% teaching: NEW-U10-VOC-TOTWAEMAN
% mdparagraph:P158:begin
The earlier \textbf{tōtwǣman}, “separate”, can now also separate beings: the person, animal or group separated is accusative, and \textbf{fram + dative} names what it is separated from. Use its familiar indicative forms. The source reading has a different verb, tōsceadan, whose exact finite form is supplied there.
% mdparagraph:P158:end

% teaching: NEW-U10-GRAM-u10-hwonne-direct-question
% teaching: NEW-U10-VOC-HWONNE
% mdparagraph:P159:begin
Invariant \textbf{hwonne}, “when?”, introduces a direct question about time. The finite verb expresses the tense and person; use the ordinary interrogative order with hwonne before it and the subject following. English auxiliary “do/\allowbreak{}did” is not an additional Old English verb in these questions. The past form of “see” itself expresses “did see”. This question word does not yet introduce a general system of temporal subordinate clauses.
% mdparagraph:P159:end

% teaching: NEW-U10-GRAM-u10-coordinated-predicate-ellipsis
% mdparagraph:P160:begin
In a coordination with a new expressed subject, a repeated predicate may be left understood from the preceding clause. The following subject still has its own nominative case. This accounts for the two coming subjects in Cook III.15.9 and for the angels' understood coming in the source reading. Supply the repeated action in understanding; this bounded pattern does not permit arbitrary deletion of verbs.
% mdparagraph:P160:end

% teaching: NEW-U10-VOC-MICEL-COMP
% mdparagraph:P161:begin
The familiar comparative \textbf{māra}, neuter \textbf{māre}, can express greater value or importance, “more than”. Use that evaluation sense in the comparison of life and raiment, with the already learned comparative connection ðonne. No new comparative formation is needed.
% mdparagraph:P161:end

\reserve{7\baselineskip}
% mdparagraph:P162:begin
\section*{9. New practice}
% mdparagraph:P162:end

\reserve{7\baselineskip}
% mdparagraph:P163:begin
\subsection*{Forms}
% mdparagraph:P163:end

\reserve{5\baselineskip}
\par\begin{minipage}{\linewidth}\raggedright
% mdparagraph:P164:begin
\textbf{F01.}
% mdparagraph:P164:end

% item: NEW-U10-F01
% mdparagraph:P165:begin
Give the complete nominative, genitive, dative and accusative singular and plural of fæder. Identify the selected endingless singular cells.
% mdparagraph:P165:end

\end{minipage}\par\addvspace{4pt}

\par\begin{minipage}{\linewidth}\raggedright
% mdparagraph:P166:begin
\textbf{F02.}
% mdparagraph:P166:end

% item: NEW-U10-F02
% mdparagraph:P167:begin
Give the complete nominative, genitive, dative and accusative singular and plural of gāt. Name every cell represented by gǣt in the selected paradigm.
% mdparagraph:P167:end

\end{minipage}\par\addvspace{4pt}

\par\begin{minipage}{\linewidth}\raggedright
% mdparagraph:P168:begin
\textbf{F03.}
% mdparagraph:P168:end

% item: NEW-U10-F03
% mdparagraph:P169:begin
Give all persons of the present and preterite indicative of forwiernan. Keep its present and preterite stems distinct.
% mdparagraph:P169:end

\end{minipage}\par\addvspace{4pt}

\par\begin{minipage}{\linewidth}\raggedright
% mdparagraph:P170:begin
\textbf{F04.}
% mdparagraph:P170:end

% item: NEW-U10-F04
% mdparagraph:P171:begin
Give the supported present and preterite indicative of unnan, marking the unsupported second-singular present cell with a dash. Give its singular and plural subjunctive forms in both tenses and identify the forms used for a grant or wish.
% mdparagraph:P171:end

\end{minipage}\par\addvspace{4pt}

\par\begin{minipage}{\linewidth}\raggedright
% mdparagraph:P172:begin
\textbf{F05.}
% mdparagraph:P172:end

% item: NEW-U10-F05
% mdparagraph:P173:begin
Give the second and third singular and third plural present indicative of fēdan, its first singular and third plural preterite, and both commands.
% mdparagraph:P173:end

\end{minipage}\par\addvspace{4pt}

\par\begin{minipage}{\linewidth}\raggedright
% mdparagraph:P174:begin
\textbf{F06.}
% mdparagraph:P174:end

% item: NEW-U10-F06
% mdparagraph:P175:begin
Give the distinct masculine/\allowbreak{}neuter singular instrumental forms of se and strong gōd. Then give the corresponding dative forms. State how the feminine singular and plural express the instrumental relation.
% mdparagraph:P175:end

\end{minipage}\par\addvspace{4pt}

\reserve{7\baselineskip}
% mdparagraph:P176:begin
\subsection*{Analysis and contrasts}
% mdparagraph:P176:end

\reserve{5\baselineskip}
\par\begin{minipage}{\linewidth}\raggedright
% mdparagraph:P177:begin
\textbf{A01.}
% mdparagraph:P177:end

% item: NEW-U10-A01
% mdparagraph:P178:begin
Translate both sentences. Identify the genitive in each and explain what governs it.
% mdparagraph:P178:end

% mdparagraph:P179:begin
\textbf{a.} Ðæs cuman sunu cōm.
% mdparagraph:P179:end

% mdparagraph:P180:begin
\textbf{b.} Ic geswāc ðæs weorces.
% mdparagraph:P180:end

\end{minipage}\par\addvspace{4pt}

\par\begin{minipage}{\linewidth}\raggedright
% mdparagraph:P181:begin
\textbf{A02.}
% mdparagraph:P181:end

% item: NEW-U10-A02
% mdparagraph:P182:begin
Translate the sentence. Identify the recipient, the thing given, and the possessor of the recipient; give their cases.
% mdparagraph:P182:end

% mdparagraph:P183:begin
Se munuc sealde ðæs cuman suna męte.
% mdparagraph:P183:end

\end{minipage}\par\addvspace{4pt}

\par\begin{minipage}{\linewidth}\raggedright
% mdparagraph:P184:begin
\textbf{A03.}
% mdparagraph:P184:end

% item: NEW-U10-A03
% mdparagraph:P185:begin
Translate both sentences and explain the difference in the phrases beginning ān.
% mdparagraph:P185:end

% mdparagraph:P186:begin
\textbf{a.} Ān munuc cōm.
% mdparagraph:P186:end

% mdparagraph:P187:begin
\textbf{b.} Ān ðāra munuca cōm.
% mdparagraph:P187:end

\end{minipage}\par\addvspace{4pt}

\par\begin{minipage}{\linewidth}\raggedright
% mdparagraph:P188:begin
\textbf{A04.}
% mdparagraph:P188:end

% item: NEW-U10-A04
% mdparagraph:P189:begin
Translate both sentences. Explain the different case of scēap and scēapa.
% mdparagraph:P189:end

% mdparagraph:P190:begin
\textbf{a.} Ic hæbbe twā scēap.
% mdparagraph:P190:end

% mdparagraph:P191:begin
\textbf{b.} Ic hæbbe fela scēapa.
% mdparagraph:P191:end

\end{minipage}\par\addvspace{4pt}

\par\begin{minipage}{\linewidth}\raggedright
% mdparagraph:P192:begin
\textbf{A05.}
% mdparagraph:P192:end

% item: NEW-U10-A05
% mdparagraph:P193:begin
Translate both sentences. Explain why the person has the same case but the food has different cases.
% mdparagraph:P193:end

% mdparagraph:P194:begin
\textbf{a.} Hē sealde ðǣm cuman męte.
% mdparagraph:P194:end

% mdparagraph:P195:begin
\textbf{b.} Hē forwiernde ðǣm cuman mętes.
% mdparagraph:P195:end

\end{minipage}\par\addvspace{4pt}

\par\begin{minipage}{\linewidth}\raggedright
% mdparagraph:P196:begin
\textbf{A06.}
% mdparagraph:P196:end

% item: NEW-U10-A06
% mdparagraph:P197:begin
Translate the sentence. Identify who is useful and to whom; explain the forms of gāt, scēaphierde and behēfu.
% mdparagraph:P197:end

% mdparagraph:P198:begin
Sēo gāt is ðǣm scēaphierde behēfu.
% mdparagraph:P198:end

\end{minipage}\par\addvspace{4pt}

\par\begin{minipage}{\linewidth}\raggedright
% mdparagraph:P199:begin
\textbf{A07.}
% mdparagraph:P199:end

% item: NEW-U10-A07
% mdparagraph:P200:begin
Translate the sentence and identify the experiencer. Explain why mē is not a nominative subject.
% mdparagraph:P200:end

% mdparagraph:P201:begin
Mē ðyncð ðis gōd.
% mdparagraph:P201:end

\end{minipage}\par\addvspace{4pt}

\par\begin{minipage}{\linewidth}\raggedright
% mdparagraph:P202:begin
\textbf{A08.}
% mdparagraph:P202:end

% item: NEW-U10-A08
% mdparagraph:P203:begin
Translate both sentences. Compare the case of the pronoun and the agreement of the finite verb.
% mdparagraph:P203:end

% mdparagraph:P204:begin
\textbf{a.} Hine hyngrede.
% mdparagraph:P204:end

% mdparagraph:P205:begin
\textbf{b.} Hē wæs hungrig.
% mdparagraph:P205:end

\end{minipage}\par\addvspace{4pt}

\par\begin{minipage}{\linewidth}\raggedright
% mdparagraph:P206:begin
\textbf{A09.}
% mdparagraph:P206:end

% item: NEW-U10-A09
% mdparagraph:P207:begin
Translate both sentences. Distinguish the relations expressed by mid and state the case it governs in each.
% mdparagraph:P207:end

% mdparagraph:P208:begin
\textbf{a.} Se munuc cōm mid ðǣm cuman.
% mdparagraph:P208:end

% mdparagraph:P209:begin
\textbf{b.} Se munuc fēdde ðone cuman mid gōdum męte.
% mdparagraph:P209:end

\end{minipage}\par\addvspace{4pt}

\par\begin{minipage}{\linewidth}\raggedright
% mdparagraph:P210:begin
\textbf{A10.}
% mdparagraph:P210:end

% item: NEW-U10-A10
% mdparagraph:P211:begin
Translate both sentences. Explain the form of gōde or ðȳ and state whether the noun ending alone distinguishes dative from instrumental.
% mdparagraph:P211:end

% mdparagraph:P212:begin
\textbf{a.} Se munuc fēdde ðone cuman gōde męte.
% mdparagraph:P212:end

% mdparagraph:P213:begin
\textbf{b.} Se munuc fēdde ðone cuman ðȳ męte.
% mdparagraph:P213:end

\end{minipage}\par\addvspace{4pt}

\par\begin{minipage}{\linewidth}\raggedright
% mdparagraph:P214:begin
\textbf{A11.}
% mdparagraph:P214:end

% item: NEW-U10-A11
% mdparagraph:P215:begin
Translate the sentence. State the case, number and gender of cuman and untrumne, and explain the agreement.
% mdparagraph:P215:end

% mdparagraph:P216:begin
Ic geseah ðone cuman untrumne.
% mdparagraph:P216:end

\end{minipage}\par\addvspace{4pt}

\par\begin{minipage}{\linewidth}\raggedright
% mdparagraph:P217:begin
\textbf{A12.}
% mdparagraph:P217:end

% item: NEW-U10-A12
% mdparagraph:P218:begin
Translate both sentences. Explain the number and function of each adjective and the difference in reference.
% mdparagraph:P218:end

% mdparagraph:P219:begin
\textbf{a.} Se munuc scrȳt ðā nacodan.
% mdparagraph:P219:end

% mdparagraph:P220:begin
\textbf{b.} Se munuc scrȳt ðone nacodan cuman.
% mdparagraph:P220:end

\end{minipage}\par\addvspace{4pt}

\reserve{7\baselineskip}
% mdparagraph:P221:begin
\subsection*{Controlled composition}
% mdparagraph:P221:end

\reserve{5\baselineskip}
\par\begin{minipage}{\linewidth}\raggedright
% mdparagraph:P222:begin
\textbf{T01.}
% mdparagraph:P222:end

% item: NEW-U10-T01
% mdparagraph:P223:begin
Translate: The monks deny the stranger food. Use forwiernan.
% mdparagraph:P223:end

\end{minipage}\par\addvspace{4pt}

\par\begin{minipage}{\linewidth}\raggedright
% mdparagraph:P224:begin
\textbf{T02.}
% mdparagraph:P224:end

% item: NEW-U10-T02
% mdparagraph:P225:begin
Translate: I grant the stranger food. Use unnan.
% mdparagraph:P225:end

\end{minipage}\par\addvspace{4pt}

\par\begin{minipage}{\linewidth}\raggedright
% mdparagraph:P226:begin
\textbf{T03.}
% mdparagraph:P226:end

% item: NEW-U10-T03
% mdparagraph:P227:begin
Translate both sentences: Some of the monks came with the stranger. Which of the monks came? Use sum and hwilc with a genitive whole.
% mdparagraph:P227:end

\end{minipage}\par\addvspace{4pt}

\par\begin{minipage}{\linewidth}\raggedright
% mdparagraph:P228:begin
\textbf{T04.}
% mdparagraph:P228:end

% item: NEW-U10-T04
% mdparagraph:P229:begin
Translate: The shepherd has many sheep. Use fela.
% mdparagraph:P229:end

\end{minipage}\par\addvspace{4pt}

\par\begin{minipage}{\linewidth}\raggedright
% mdparagraph:P230:begin
\textbf{T05.}
% mdparagraph:P230:end

% item: NEW-U10-T05
% mdparagraph:P231:begin
Translate: The monks feed the strangers with good food. Express the means without a preposition.
% mdparagraph:P231:end

\end{minipage}\par\addvspace{4pt}

\par\begin{minipage}{\linewidth}\raggedright
% mdparagraph:P232:begin
\textbf{T06.}
% mdparagraph:P232:end

% item: NEW-U10-T06
% mdparagraph:P233:begin
Translate: With that food I fed the sick. Use the distinct instrumental of se and an adjective naming the sick people.
% mdparagraph:P233:end

\end{minipage}\par\addvspace{4pt}

\par\begin{minipage}{\linewidth}\raggedright
% mdparagraph:P234:begin
\textbf{T07.}
% mdparagraph:P234:end

% item: NEW-U10-T07
% mdparagraph:P235:begin
Translate: They were hungry. Use the impersonal verb hyngran.
% mdparagraph:P235:end

\end{minipage}\par\addvspace{4pt}

\par\begin{minipage}{\linewidth}\raggedright
% mdparagraph:P236:begin
\textbf{T08.}
% mdparagraph:P236:end

% item: NEW-U10-T08
% mdparagraph:P237:begin
Translate: That seemed good to us. Use the chapter’s ðyncan construction.
% mdparagraph:P237:end

\end{minipage}\par\addvspace{4pt}

\reserve{7\baselineskip}
% mdparagraph:P238:begin
\section*{10. Cook III.13–15}
% mdparagraph:P238:end

% mdparagraph:P239:begin
Translate into Old English. Preserve the stated number of an addressee and the number of a named object. Where a prompt leaves a second-person number unstated, choose a consistent reading. “Om.” instructs omission of the specified English word; it is not vocabulary to translate. Direct speech stays direct speech. The English “will” and “shall” should be interpreted by the meaning, as explained in §8.
% mdparagraph:P239:end

% mdparagraph:P240:begin
The complete original sets follow. Their grammatical references are preserved. The two local helps supply expressions beyond the independent composition scope; use them for their stated items.
% mdparagraph:P240:end

\reserve{7\baselineskip}
% mdparagraph:P241:begin
\subsection*{III.13}
% mdparagraph:P241:end

\reserve{5\baselineskip}
\par\begin{minipage}{\linewidth}\raggedright
% mdparagraph:P242:begin
\textbf{III.13.01.}
% mdparagraph:P242:end

% item: EX-III-13-01
% mdparagraph:P243:begin
Christ is the Son of God [God’s Son].
% mdparagraph:P243:end

\end{minipage}\par\addvspace{4pt}

\par\begin{minipage}{\linewidth}\raggedright
% mdparagraph:P244:begin
\textbf{III.13.02.}
% mdparagraph:P244:end

% item: EX-III-13-02
% mdparagraph:P245:begin
Christ is our Lord.
% mdparagraph:P245:end

\end{minipage}\par\addvspace{4pt}

\par\begin{minipage}{\linewidth}\raggedright
% mdparagraph:P246:begin
\textbf{III.13.03.}
% mdparagraph:P246:end

% item: EX-III-13-03
% mdparagraph:P247:begin
When saw I you (sing.) hungry?
% mdparagraph:P247:end

\end{minipage}\par\addvspace{4pt}

\par\begin{minipage}{\linewidth}\raggedright
% mdparagraph:P248:begin
\textbf{III.13.04.}
% mdparagraph:P248:end

% item: EX-III-13-04
% mdparagraph:P249:begin
Gold is precious.
% mdparagraph:P249:end

\end{minipage}\par\addvspace{4pt}

\par\begin{minipage}{\linewidth}\raggedright
% mdparagraph:P250:begin
\textbf{III.13.05.}
% mdparagraph:P250:end

% item: EX-III-13-05
% mdparagraph:P251:begin
The Day of Judgment will come.
% mdparagraph:P251:end

\end{minipage}\par\addvspace{4pt}

\par\begin{minipage}{\linewidth}\raggedright
% mdparagraph:P252:begin
\textbf{III.13.06.}
% mdparagraph:P252:end

% item: EX-III-13-06
% mdparagraph:P253:begin
The farmer says to the workman, I will sell this house.
% mdparagraph:P253:end

\end{minipage}\par\addvspace{4pt}

\par\begin{minipage}{\linewidth}\raggedright
% mdparagraph:P254:begin
\textbf{III.13.07.}
% mdparagraph:P254:end

% item: EX-III-13-07
% mdparagraph:P255:begin
The workman answers the farmer, I will buy your house.
% mdparagraph:P255:end

\end{minipage}\par\addvspace{4pt}

\par\begin{minipage}{\linewidth}\raggedright
% mdparagraph:P256:begin
\textbf{III.13.08.}
% mdparagraph:P256:end

% item: EX-III-13-08
% mdparagraph:P257:begin
Life is more than raiment.
% mdparagraph:P257:end

\end{minipage}\par\addvspace{4pt}

\par\begin{minipage}{\linewidth}\raggedright
% mdparagraph:P258:begin
\textbf{III.13.09.}
% mdparagraph:P258:end

% item: EX-III-13-09
% mdparagraph:P259:begin
The stranger is in prison.
% mdparagraph:P259:end

\end{minipage}\par\addvspace{4pt}

\par\begin{minipage}{\linewidth}\raggedright
% mdparagraph:P260:begin
\textbf{III.13.10.}
% mdparagraph:P260:end

% item: EX-III-13-10
% mdparagraph:P261:begin
The merchant is in the inn.
% mdparagraph:P261:end

\end{minipage}\par\addvspace{4pt}

\par\begin{minipage}{\linewidth}\raggedright
% mdparagraph:P262:begin
\textbf{III.13.11.}
% mdparagraph:P262:end

% item: EX-III-13-11
% mdparagraph:P263:begin
The father clothes the son.
% mdparagraph:P263:end

\end{minipage}\par\addvspace{4pt}

\par\begin{minipage}{\linewidth}\raggedright
% mdparagraph:P264:begin
\textbf{III.13.12.}
% mdparagraph:P264:end

% item: EX-III-13-12
% mdparagraph:P265:begin
This world is great.
% mdparagraph:P265:end

\end{minipage}\par\addvspace{4pt}

\par\begin{minipage}{\linewidth}\raggedright
% mdparagraph:P266:begin
\textbf{III.13.13.}
% mdparagraph:P266:end

% item: EX-III-13-13
% mdparagraph:P267:begin
The shepherd feeds the sheep (sing.).
% mdparagraph:P267:end

\end{minipage}\par\addvspace{4pt}

\par\begin{minipage}{\linewidth}\raggedright
% mdparagraph:P268:begin
\textbf{III.13.14.}
% mdparagraph:P268:end

% item: EX-III-13-14
% mdparagraph:P269:begin
The merchants are rich.
% mdparagraph:P269:end

\end{minipage}\par\addvspace{4pt}

\par\begin{minipage}{\linewidth}\raggedright
% mdparagraph:P270:begin
\textbf{III.13.15.}
% mdparagraph:P270:end

% item: EX-III-13-15
% mdparagraph:P271:begin
The soldier drinks wine.
% mdparagraph:P271:end

\end{minipage}\par\addvspace{4pt}

\par\begin{minipage}{\linewidth}\raggedright
% mdparagraph:P272:begin
\textbf{III.13.16.}
% mdparagraph:P272:end

% item: EX-III-13-16
% mdparagraph:P273:begin
The monks drank water.
% mdparagraph:P273:end

\end{minipage}\par\addvspace{4pt}

\reserve{7\baselineskip}
% mdparagraph:P274:begin
\subsection*{III.14}
% mdparagraph:P274:end

\reserve{5\baselineskip}
\par\begin{minipage}{\linewidth}\raggedright
% mdparagraph:P275:begin
\textbf{III.14.01.}
% mdparagraph:P275:end

% item: EX-III-14-01
% mdparagraph:P276:begin
The king sat on the throne.
% mdparagraph:P276:end

\end{minipage}\par\addvspace{4pt}

\par\begin{minipage}{\linewidth}\raggedright
% mdparagraph:P277:begin
\textbf{III.14.02.}
% mdparagraph:P277:end

% item: EX-III-14-02
% mdparagraph:P278:begin
Stand on my right hand.
% mdparagraph:P278:end

\end{minipage}\par\addvspace{4pt}

\par\begin{minipage}{\linewidth}\raggedright
% mdparagraph:P279:begin
\textbf{III.14.03.}
% mdparagraph:P279:end

% item: EX-III-14-03
% mdparagraph:P280:begin
All nations shall come to the judgment.
% mdparagraph:P280:end

\end{minipage}\par\addvspace{4pt}

\par\begin{minipage}{\linewidth}\raggedright
% mdparagraph:P281:begin
\textbf{III.14.04.}
% mdparagraph:P281:end

% item: EX-III-14-04
% mdparagraph:P282:begin
The devil is (om. the) father of the unrighteous.
% mdparagraph:P282:end

\end{minipage}\par\addvspace{4pt}

\par\begin{minipage}{\linewidth}\raggedright
% mdparagraph:P283:begin
\textbf{III.14.05.}
% mdparagraph:P283:end

% item: EX-III-14-05
% mdparagraph:P284:begin
All men are sinful.
% mdparagraph:P284:end

\end{minipage}\par\addvspace{4pt}

\par\begin{minipage}{\linewidth}\raggedright
% mdparagraph:P285:begin
\textbf{III.14.06.}
% mdparagraph:P285:end

% item: EX-III-14-06
% mdparagraph:P286:begin
The priest is sick.
% mdparagraph:P286:end

\end{minipage}\par\addvspace{4pt}

\par\begin{minipage}{\linewidth}\raggedright
% mdparagraph:P287:begin
\textbf{III.14.07.}
% mdparagraph:P287:end

% item: EX-III-14-07
% mdparagraph:P288:begin
Grant (sing.) me garments.
% mdparagraph:P288:end

\end{minipage}\par\addvspace{4pt}

\par\begin{minipage}{\linewidth}\raggedright
% mdparagraph:P289:begin
\textbf{III.14.08.}
% mdparagraph:P289:end

% item: EX-III-14-08
% mdparagraph:P290:begin
The horse is useful to the king.
% mdparagraph:P290:end

\end{minipage}\par\addvspace{4pt}

\par\begin{minipage}{\linewidth}\raggedright
% mdparagraph:P291:begin
\textbf{III.14.09.}
% mdparagraph:P291:end

% item: EX-III-14-09
% mdparagraph:P292:begin
The loss of all your things was strange.
% mdparagraph:P292:end

\end{minipage}\par\addvspace{4pt}

\par\begin{minipage}{\linewidth}\raggedright
% mdparagraph:P293:begin
\textbf{III.14.10.}
% mdparagraph:P293:end

% item: EX-III-14-10
% mdparagraph:P294:begin
He departed from her.
% mdparagraph:P294:end

\end{minipage}\par\addvspace{4pt}

\par\begin{minipage}{\linewidth}\raggedright
% mdparagraph:P295:begin
\textbf{III.14.11.}
% mdparagraph:P295:end

% item: EX-III-14-11
% mdparagraph:P296:begin
Give (sing.) me to drink.
% mdparagraph:P296:end

\end{minipage}\par\addvspace{4pt}

% teaching: NEW-U10-COOK-HELPS
% teaching: NEW-U10-GRAM-u10-sellan-drink-infinitive
% support: U10-EXHELP-DRINK
% mdparagraph:P297:begin
\textbf{Local help.} For give someone to drink, use sęllan with the person in the dative and plain drincan, without tō. Apply the already known singular command and pronoun forms. This entire construction is supplied for this item.
% mdparagraph:P297:end

\reserve{5\baselineskip}
\par\begin{minipage}{\linewidth}\raggedright
% mdparagraph:P298:begin
\textbf{III.14.12.}
% mdparagraph:P298:end

% item: EX-III-14-12
% mdparagraph:P299:begin
Visit (plur.) the sick.
% mdparagraph:P299:end

\end{minipage}\par\addvspace{4pt}

\par\begin{minipage}{\linewidth}\raggedright
% mdparagraph:P300:begin
\textbf{III.14.13.}
% mdparagraph:P300:end

% item: EX-III-14-13
% mdparagraph:P301:begin
When did you (sing.) see the goat?
% mdparagraph:P301:end

\end{minipage}\par\addvspace{4pt}

\par\begin{minipage}{\linewidth}\raggedright
% mdparagraph:P302:begin
\textbf{III.14.14.}
% mdparagraph:P302:end

% item: EX-III-14-14
% mdparagraph:P303:begin
When was he in the inn?
% mdparagraph:P303:end

\end{minipage}\par\addvspace{4pt}

\par\begin{minipage}{\linewidth}\raggedright
% mdparagraph:P304:begin
\textbf{III.14.15.}
% mdparagraph:P304:end

% item: EX-III-14-15
% mdparagraph:P305:begin
This is eternal life.
% mdparagraph:P305:end

\end{minipage}\par\addvspace{4pt}

\par\begin{minipage}{\linewidth}\raggedright
% mdparagraph:P306:begin
\textbf{III.14.16.}
% mdparagraph:P306:end

% item: EX-III-14-16
% mdparagraph:P307:begin
The prisons were empty.
% mdparagraph:P307:end

\end{minipage}\par\addvspace{4pt}

\par\begin{minipage}{\linewidth}\raggedright
% mdparagraph:P308:begin
\textbf{III.14.17.}
% mdparagraph:P308:end

% item: EX-III-14-17
% mdparagraph:P309:begin
The living God is eternal.
% mdparagraph:P309:end

\end{minipage}\par\addvspace{4pt}

\reserve{7\baselineskip}
% mdparagraph:P310:begin
\subsection*{III.15}
% mdparagraph:P310:end

\reserve{5\baselineskip}
\par\begin{minipage}{\linewidth}\raggedright
% mdparagraph:P311:begin
\textbf{III.15.01.}
% mdparagraph:P311:end

% item: EX-III-15-01
% mdparagraph:P312:begin
He stood before the king.
% mdparagraph:P312:end

\end{minipage}\par\addvspace{4pt}

\par\begin{minipage}{\linewidth}\raggedright
% mdparagraph:P313:begin
\textbf{III.15.02.}
% mdparagraph:P313:end

% item: EX-III-15-02
% mdparagraph:P314:begin
The goat eats grass.
% mdparagraph:P314:end

\end{minipage}\par\addvspace{4pt}

\par\begin{minipage}{\linewidth}\raggedright
% mdparagraph:P315:begin
\textbf{III.15.03.}
% mdparagraph:P315:end

% item: EX-III-15-03
% mdparagraph:P316:begin
The angel flies to the earth.
% mdparagraph:P316:end

\end{minipage}\par\addvspace{4pt}

\par\begin{minipage}{\linewidth}\raggedright
% mdparagraph:P317:begin
\textbf{III.15.04.}
% mdparagraph:P317:end

% item: EX-III-15-04
% mdparagraph:P318:begin
You (sing.) deny me food.
% mdparagraph:P318:end

\end{minipage}\par\addvspace{4pt}

\par\begin{minipage}{\linewidth}\raggedright
% mdparagraph:P319:begin
\textbf{III.15.05.}
% mdparagraph:P319:end

% item: EX-III-15-05
% mdparagraph:P320:begin
The monk did it in Christ’s name.
% mdparagraph:P320:end

\end{minipage}\par\addvspace{4pt}

\par\begin{minipage}{\linewidth}\raggedright
% mdparagraph:P321:begin
\textbf{III.15.06.}
% mdparagraph:P321:end

% item: EX-III-15-06
% mdparagraph:P322:begin
He will separate the goats from the sheep.
% mdparagraph:P322:end

\end{minipage}\par\addvspace{4pt}

\par\begin{minipage}{\linewidth}\raggedright
% mdparagraph:P323:begin
\textbf{III.15.07.}
% mdparagraph:P323:end

% item: EX-III-15-07
% mdparagraph:P324:begin
He poured her out wine and water.
% mdparagraph:P324:end

\end{minipage}\par\addvspace{4pt}

\par\begin{minipage}{\linewidth}\raggedright
% mdparagraph:P325:begin
\textbf{III.15.08.}
% mdparagraph:P325:end

% item: EX-III-15-08
% mdparagraph:P326:begin
The Lord will come with his angels.
% mdparagraph:P326:end

\end{minipage}\par\addvspace{4pt}

\par\begin{minipage}{\linewidth}\raggedright
% mdparagraph:P327:begin
\textbf{III.15.09.}
% mdparagraph:P327:end

% item: EX-III-15-09
% mdparagraph:P328:begin
The morning cometh, and also the night.
% mdparagraph:P328:end

\end{minipage}\par\addvspace{4pt}

\par\begin{minipage}{\linewidth}\raggedright
% mdparagraph:P329:begin
\textbf{III.15.10.}
% mdparagraph:P329:end

% item: EX-III-15-10
% mdparagraph:P330:begin
Work the works of God.
% mdparagraph:P330:end

\end{minipage}\par\addvspace{4pt}

\par\begin{minipage}{\linewidth}\raggedright
% mdparagraph:P331:begin
\textbf{III.15.11.}
% mdparagraph:P331:end

% item: EX-III-15-11
% mdparagraph:P332:begin
Clothe (sing.) the naked.
% mdparagraph:P332:end

\end{minipage}\par\addvspace{4pt}

\par\begin{minipage}{\linewidth}\raggedright
% mdparagraph:P333:begin
\textbf{III.15.12.}
% mdparagraph:P333:end

% item: EX-III-15-12
% mdparagraph:P334:begin
I will not sit on the king’s throne.
% mdparagraph:P334:end

\end{minipage}\par\addvspace{4pt}

\par\begin{minipage}{\linewidth}\raggedright
% mdparagraph:P335:begin
\textbf{III.15.13.}
% mdparagraph:P335:end

% item: EX-III-15-13
% mdparagraph:P336:begin
The bird sees the light.
% mdparagraph:P336:end

\end{minipage}\par\addvspace{4pt}

\par\begin{minipage}{\linewidth}\raggedright
% mdparagraph:P337:begin
\textbf{III.15.14.}
% mdparagraph:P337:end

% item: EX-III-15-14
% mdparagraph:P338:begin
The men will see the stars.
% mdparagraph:P338:end

\end{minipage}\par\addvspace{4pt}

\par\begin{minipage}{\linewidth}\raggedright
% mdparagraph:P339:begin
\textbf{III.15.15.}
% mdparagraph:P339:end

% item: EX-III-15-15
% mdparagraph:P340:begin
As long as you (plur.) fed the hungry, you did it to the Lord.
% mdparagraph:P340:end

\end{minipage}\par\addvspace{4pt}

% teaching: NEW-U10-GRAM-u10-swa-lange-swa-supplied
% support: U10-EXHELP-AS-LONG-AS
% mdparagraph:P341:begin
\textbf{Local help.} Use the supplied connection swā lange swā for the opening As long as. Cook’s note17 explains the biblical expression as inasmuch as: the deeds done to the hungry are counted as done to the Lord. Both swā words belong to this whole connection; the earlier so/\allowbreak{}thus sense does not teach it. The two you references are plural; use the taught past forms and case frames.
% mdparagraph:P341:end

\reserve{5\baselineskip}
\par\begin{minipage}{\linewidth}\raggedright
% mdparagraph:P342:begin
\textbf{III.15.16.}
% mdparagraph:P342:end

% item: EX-III-15-16
% mdparagraph:P343:begin
When did you see the stranger sick?
% mdparagraph:P343:end

\end{minipage}\par\addvspace{4pt}

\par\begin{minipage}{\linewidth}\raggedright
% mdparagraph:P344:begin
\textbf{III.15.17.}
% mdparagraph:P344:end

% item: EX-III-15-17
% mdparagraph:P345:begin
The hungry man came to the ships.
% mdparagraph:P345:end

\end{minipage}\par\addvspace{4pt}

\par\begin{minipage}{\linewidth}\raggedright
% mdparagraph:P346:begin
\textbf{III.15.18.}
% mdparagraph:P346:end

% item: EX-III-15-18
% mdparagraph:P347:begin
The soldier brought us the horses’ fodder.
% mdparagraph:P347:end

\end{minipage}\par\addvspace{4pt}

\par\begin{minipage}{\linewidth}\raggedright
% mdparagraph:P348:begin
\textbf{III.15.19.}
% mdparagraph:P348:end

% item: EX-III-15-19
% mdparagraph:P349:begin
The farmer has the blacksmith’s hammer.
% mdparagraph:P349:end

\end{minipage}\par\addvspace{4pt}

\par\begin{minipage}{\linewidth}\raggedright
% mdparagraph:P350:begin
\textbf{III.15.20.}
% mdparagraph:P350:end

% item: EX-III-15-20
% mdparagraph:P351:begin
Who has the tailor’s needles?
% mdparagraph:P351:end

\end{minipage}\par\addvspace{4pt}

\par\begin{minipage}{\linewidth}\raggedright
% mdparagraph:P352:begin
\textbf{III.15.21.}
% mdparagraph:P352:end

% item: EX-III-15-21
% mdparagraph:P353:begin
I hear the bellows (169) blow.
% mdparagraph:P353:end

\end{minipage}\par\addvspace{4pt}

\par\begin{minipage}{\linewidth}\raggedright
% mdparagraph:P354:begin
\textbf{III.15.22.}
% mdparagraph:P354:end

% item: EX-III-15-22
% mdparagraph:P355:begin
Do (om.) you (sing.) see the hammers smite?
% mdparagraph:P355:end

% mdparagraph:P356:begin
In III.15.22, \textbf{(om.) applies to English Do}. Use the finite interrogative construction with “see” and its accusative actor plus infinitive; do not add a translation of auxiliary do. In III.15.21, (169) is Cook's retained reference to his perception construction.
% mdparagraph:P356:end

\end{minipage}\par\addvspace{4pt}

\reserve{7\baselineskip}
% mdparagraph:P357:begin
\section*{11. The stranger and the monk}
% mdparagraph:P357:end

% mdparagraph:P358:begin
New teaching material. This connected account uses the vocabulary and constructions taught above.
% mdparagraph:P358:end

% composition: NEW-U10-READING-01
% mdparagraph:P359:begin
\par\begin{minipage}{\linewidth}\raggedright
\readingtext{Sum cuma cōm tō ðǣm giesthūse. Hine hyngrede, and hē wæs untrum. Ān ðāra munuca geseah ðone cuman. Se munuc sealde him męte and wæter. Se cuma cwæð: “Ðū ne forwiernst mē mętes.” Se munuc cwæð: “Ic fēde ðē gōde męte.” Se cuma ǣt ðone męte and dranc ðæt wæter.}
\end{minipage}\par
% mdparagraph:P359:end

\reserve{5\baselineskip}
\par\begin{minipage}{\linewidth}\raggedright
% mdparagraph:P360:begin
\textbf{R01.}
% mdparagraph:P360:end

% item: NEW-U10-R01
% mdparagraph:P361:begin
In the new account, identify the referents of hine, him, mē and ðē at each occurrence. State their cases and the words or constructions governing them.
% mdparagraph:P361:end

\end{minipage}\par\addvspace{4pt}

\par\begin{minipage}{\linewidth}\raggedright
% mdparagraph:P362:begin
\textbf{R02.}
% mdparagraph:P362:end

% item: NEW-U10-R02
% mdparagraph:P363:begin
In the new account, explain the case and function of ðāra munuca and mętes. Distinguish the functions of the forms cuma and cuman.
% mdparagraph:P363:end

\end{minipage}\par\addvspace{4pt}

\par\begin{minipage}{\linewidth}\raggedright
% mdparagraph:P364:begin
\textbf{R03.}
% mdparagraph:P364:end

% item: NEW-U10-R03
% mdparagraph:P365:begin
In the monk’s final reply, explain the ending of gōde. Rewrite only that means phrase using mid.
% mdparagraph:P365:end

\end{minipage}\par\addvspace{4pt}

\par\begin{minipage}{\linewidth}\raggedright
% mdparagraph:P366:begin
\textbf{R04.}
% mdparagraph:P366:end

% item: NEW-U10-R04
% mdparagraph:P367:begin
Translate the whole new account. Preserve who gives, who is hungry and who eats or drinks.
% mdparagraph:P367:end

\end{minipage}\par\addvspace{4pt}

\reserve{7\baselineskip}
% mdparagraph:P368:begin
\section*{12. The Day of Judgment}
% mdparagraph:P368:end

% mdparagraph:P369:begin
Exact source passage. The three complete opening sentences form a scene of arrival, enthronement, gathering and separation. The passage ends before the king's speech. Its original spellings and quantities are retained. The supplied help belongs to this reading; expressions marked there are not additional unaided composition vocabulary.
% mdparagraph:P369:end

\par\begin{minipage}{\linewidth}\raggedright
% source: FB-III-JUDGMENT-TABLEAU-01
% mdparagraph:P370:begin
\readingtext{Witodlice Mannes Bearn cymð on his mægenðrymme, and ealle ęnglas samod mid him tō ðǣm miclan dōme. Ðonne sitt hē on ðǣm setle his mægenðrymnesse, and bēoð gegaderode ætforan him ealle ðēoda, and hē tōscǣt hīe on twā, swā-swā scēaphierde tōscǣt scēap fram gātum. Ðonne gelōgað hē ðā scēap on his swīðran hand, and ðā gǣt on his winstran.}
% mdparagraph:P370:end

\end{minipage}\par
\reserve{7\baselineskip}
% mdparagraph:P371:begin
\subsection*{Local help}
% mdparagraph:P371:end

% teaching: NEW-U10-SOURCE-HELPS
% teaching: NEW-U10-VOC-BEARN
% teaching: NEW-U10-VOC-MAEGENTHRYMM
% teaching: NEW-U10-VOC-MAEGENTHRYMNES
% teaching: NEW-U10-VOC-WINSTRE
% teaching: NEW-U10-VOC-WITODLICE
% teaching: NEW-U10-VOC-SAMOD
% teaching: NEW-U10-VOC-TOSCEADAN
% teaching: NEW-U10-VOC-GELOGIAN
% teaching: NEW-U10-VOC-AETFORAN
% teaching: NEW-U10-VOC-THONNE
% teaching: NEW-U10-GRAM-u10-genitive-quality-source
% mdparagraph:P372:begin
The nine groups below supply the difficult words and larger connections. In particular, \textbf{mægenðrymm} and \textbf{mægenðrymnes} both mean “glory, majesty”; \textbf{tōsceadan} means “separate, divide”, and \textbf{gelōgian} “place, set”. Only their actual forms in this passage are supplied. The nominal and adjective tables above do not imply complete paradigms for these source-only words.
% mdparagraph:P372:end

% mdparagraph:P373:begin
The genitive of quality characterizes a noun by an associated quality: the throne “of his majesty” is a majestic throne. This reading supplies that abstract expression. It does not require free composition with abstract genitives of quality.
% mdparagraph:P373:end

% support: U10-HELP01
% mdparagraph:P374:begin
\textbf{1. Witodlice Mannes Bearn.} Witodlice is a connective here, meaning “now; and” (Cook gives Latin autem). Mannes Bearn means the Son of Man: mannes is genitive singular and Bearn neuter nominative singular. The subsequent hē/\allowbreak{}him/\allowbreak{}his refer to this male person.
% mdparagraph:P374:end

% support: U10-HELP02
% mdparagraph:P375:begin
\textbf{2. on his mægenðrymme.} In his glory/\allowbreak{}majesty. Mægenðrymme is dative masculine singular. This abstract expression is supplied as a whole.
% mdparagraph:P375:end

% support: U10-HELP03
% mdparagraph:P376:begin
\textbf{3. and ealle ęnglas samod mid him ….} And all the angels together with him. The source leaves the coming verb understood in this coordinated part; the angels are nominative plural, not objects of with. Mid him separately means with him. Tō ðǣm miclan dōme means to the great judgment.
% mdparagraph:P376:end

% support: U10-HELP04
% mdparagraph:P377:begin
\textbf{4. Ðonne.} Then. Both occurrences are independent sentence adverbs sequencing the scene; they are not the previously taught than or subordinate when.
% mdparagraph:P377:end

% support: U10-HELP05
% mdparagraph:P378:begin
\textbf{5. setle his mægenðrymnesse.} Throne/\allowbreak{}seat of his majesty. Mægenðrymnesse is a feminine genitive singular characterizing the throne; Cook note7 refers to the genitive of quality (§153.f).
% mdparagraph:P378:end

% support: U10-HELP06
% mdparagraph:P379:begin
\textbf{6. bēoð gegaderode ætforan him ealle ðēoda.} All nations will be gathered before him. Bēoð is third-person plural, and ðēoda is the following nominative feminine plural subject. Gegaderode is the exact supplied plural passive participle; ætforan him means before him. Ealle is the feminine plural -e variant here; the ordinary selected target ealla remains valid.
% mdparagraph:P379:end

% support: U10-HELP07
% mdparagraph:P380:begin
\textbf{7. tōscǣt hīe on twā.} Separates them into two groups. Tōscǣt is third-person singular present; hīe is its accusative plural object. On twā is the locally supplied divided-in-two expression; twā is neuter accusative plural here.
% mdparagraph:P380:end

% support: U10-HELP08
% mdparagraph:P381:begin
\textbf{8. swā-swā scēaphierde tōscǣt scēap fram gātum.} Just as a shepherd separates sheep from goats. The whole comparative connection is supplied. Scēaphierde is the nominative subject, scēap the accusative plural object, and gātum dative after fram. The noun scēap has identical singular/\allowbreak{}plural nominative/\allowbreak{}accusative shapes; the source explicitly identifies this occurrence as plural.
% mdparagraph:P381:end

% support: U10-HELP09
% mdparagraph:P382:begin
\textbf{9. gelōgað … on his swīðran hand … on his winstran.} He places the sheep at his right hand and the goats at his left. Gelōgað is an exact supplied third-singular present. The two animal phrases are accusative plural objects. The right/\allowbreak{}left expressions share a placement frame; on takes accusative hand in this bounded right/\allowbreak{}left-side expression (the final hand is understood). The following unassigned source line8 uses the same expression with static standað; this is not a universal motion-versus-rest rule. Swīðran/\allowbreak{}winstran are weak feminine singular forms. Their -an alone does not distinguish accusative from dative; judge the complete clause and the bare hand form.
% mdparagraph:P382:end

% mdparagraph:P383:begin
The familiar \textbf{cymð} and \textbf{sitt} are the selected third-singular forms of cuman and sittan. \textbf{Miclan} is the weak form of micel with its reduced stem. The passage has six overt finite predicates. The angels' coming is understood in the coordination; recover it from cymð rather than treating the angels as the object of mid.
% mdparagraph:P383:end

\reserve{7\baselineskip}
% mdparagraph:P384:begin
\subsection*{Cook's original notes}
% mdparagraph:P384:end

% mdparagraph:P385:begin
The numbers and anchors below preserve Cook's apparatus separately from the passage and the local help. His historical and etymological questions are source notes, not additional exercises assigned here.
% mdparagraph:P385:end

% source-note: U10-FB134-N01
% source-note: U10-FB134-N02
% source-note: U10-FB134-N03
% source-note: U10-FB134-N04
% source-note: U10-FB134-N05
% source-note: U10-FB134-N06
% source-note: U10-FB134-N07
% source-note: U10-FB134-N08
% source-note: U10-FB134-N09
% source-note: U10-FB134-N10
% source-note: U10-FB134-N11
% source-note: U10-FB134-N12
% source-note: U10-FB134-N13
% source-note: U10-FB134-N14
% mdtable:T39:begin
\par\begingroup
\fontsize{10.5}{13.2}\selectfont
\setlength{\LTleft}{0pt}\setlength{\LTright}{\fill}
\noindent
\begin{tabular}{@{}>{\raggedright\arraybackslash}p{8mm}>{\raggedright\arraybackslash}p{31mm}>{\raggedright\arraybackslash}p{79.37mm}@{}}
\toprule
\textbf{Note} & \textbf{Anchor} & \textbf{Original text} \\
\midrule
1 & Witodlice & Lat. autem. \\
2 & cymð & See cuman. \\
3 & ęnglas & See ęngel. What is the history of this word before it entered Old English? \\
4 & miclan & See 55. \\
5 & dōme & In what modern compound does this meaning of dōm persist? \\
6 & sitt & See sittan. \\
7 & mægenðrymnesse & See 153. f. \\
8 & ðēoda & Nom. plur. \\
9 & tōscǣt & See tōsceadan. Account for the vowel ǣ. \\
10 & scēaphierde & In compound words, the vowel of the first syllable is apt to be shortened in Mod. Eng., the more general principle being that shortening is apt to occur before an accumulation of consonants. Besides scēaphierde, shepherd, note e.g. wīsdōm, wisdom. \\
11 & scēap & Plural; account for the form. \\
12 & gātum & See 24. \\
13 & swīðran & See swīð. \\
14 & gǣt & See 52. \\
\bottomrule
\end{tabular}
\par\endgroup

% mdtable:T39:end
\reserve{7\baselineskip}
% mdparagraph:P386:begin
\subsection*{Reading tasks}
% mdparagraph:P386:end

% mdparagraph:P387:begin
Use all the printed local help in these tasks.
% mdparagraph:P387:end

\reserve{5\baselineskip}
\par\begin{minipage}{\linewidth}\raggedright
% mdparagraph:P388:begin
\textbf{S01.}
% mdparagraph:P388:end

% item: NEW-U10-S01
% mdparagraph:P389:begin
Using the printed local help, identify who comes, who is gathered, and what separation follows. Give the boundaries of the corresponding clauses.
% mdparagraph:P389:end

\end{minipage}\par\addvspace{4pt}

\par\begin{minipage}{\linewidth}\raggedright
% mdparagraph:P390:begin
\textbf{S02.}
% mdparagraph:P390:end

% item: NEW-U10-S02
% mdparagraph:P391:begin
In the separation and placement clauses, identify the cases and roles of hīe, ðā scēap and ðā gǣt. Distinguish the genders of the animal nouns.
% mdparagraph:P391:end

\end{minipage}\par\addvspace{4pt}

\par\begin{minipage}{\linewidth}\raggedright
% mdparagraph:P392:begin
\textbf{S03.}
% mdparagraph:P392:end

% item: NEW-U10-S03
% mdparagraph:P393:begin
Using the printed help, explain the two occurrences of ðonne and the relation introduced by swā-swā. Distinguish a sequencing adverb from a clause-introducing expression.
% mdparagraph:P393:end

\end{minipage}\par\addvspace{4pt}

\par\begin{minipage}{\linewidth}\raggedright
% mdparagraph:P394:begin
\textbf{S04.}
% mdparagraph:P394:end

% item: NEW-U10-S04
% mdparagraph:P395:begin
Translate the whole selected Judgment passage, using the printed source notes and local help.
% mdparagraph:P395:end

\end{minipage}\par\addvspace{4pt}

\clearpage
% mdparagraph:P396:begin
\section*{Source and editorial notes}
% mdparagraph:P396:end

\begingroup\fontsize{10}{12.8}\selectfont
\setlength{\parskip}{1.5pt}
% mdparagraph:P397:begin
\itempara{1}{\textbf{Editions, locators and fidelity.} Cook, \emph{Exercises in Old English} (1895), printed pp. 6–7, one-based PDF pp. 14–15, III.13–15: all 16, 17 and 22 items respectively are retained. Line endings are joined, typographic apostrophes are represented consistently, and the space before the question mark in III.15.20 is closed: source “needles ?”, learner “needles?”. These are declared formatting normalizations, not changes to the wording. “Tailor's”, the bracketed “[God's Son]”, the explicit singular sheep and plural visiting command, the (169) reference, and the omission directions retain their source force. Cook's \emph{A First Book in Old English}, second edition (1895), is abbreviated \textbf{FB} below. Its printed page number + 16 gives the one-based PDF page; the Exercises offset is + 8.}
% mdparagraph:P397:end

% mdparagraph:P398:begin
The reading is FB Reader III, “The Day of Judgment”, printed p. 134 /\allowbreak{} PDF p. 150. Cook's attribution is: “From Ælfric’s Homilies, vol. 2, pp. 106–108; being a paraphrase of Matt. 25:31–46.” The assigned 59-word extract is its first three complete sentences. Editorial line counting gives lines 1–7, anchored by the printed marginal 5; only 5 is a printed line number in this extract. The final assigned words occur partway through editorial line 7; the following royal speech is unassigned. Notes 1–14 are reproduced above without turning their questions into assignments. Note cues and printed line wrapping are kept in the separate apparatus; no word of the passage is modernized. Following material on pp. 134–136 corroborates selected expressions and local helps, without becoming additional assigned reading.
% mdparagraph:P398:end

% mdparagraph:P399:begin
\itempara{2}{\textbf{Case relations and C02.} Dependent genitives: FB §§152–154, pp. 89–90; governed genitives with dative persons: §§155–159, pp. 91–92. Dative constructions and prepositions: §§157–167, pp. 92–95. In §153.d, p. 89 /\allowbreak{} PDF p. 105, Cook prints \textbf{in Frēan ęgesan}, “in fear of the Lord”, but says \textbf{“ęgesan is an obj. gen.”} The dependent objective genitive is \textbf{Frēan}: the Lord is the object of the fear, as Cook's own verbal paraphrase makes clear. Correction \textbf{C02} corrects this analysis and preserves the phrase. The entire citation is supplied here; its words introduce no further production vocabulary. The source quality genitive is referred to §153.f by p. 134 note 7.}
% mdparagraph:P399:end

% mdparagraph:P400:begin
\itempara{3}{\textbf{Noun models.} FB §§43–44, pp. 28–30, strong masculine models; §43.8, p. 30, exceptional fæder; §47, pp. 32–33, strong neuters; §53, p. 37, weak nouns; §§51–52, pp. 35–37, strong feminine variants and root nouns. Hand's glossary entry, p. 297, explicitly refers to §51.1.3: this licenses its endingless accusative as well as its -a oblique variant. Woruld is named among the endingless-accusative examples. Full-headword/\allowbreak{}short-stem behavior is checked against each lexical entry; it is not imposed on compounds automatically.}
% mdparagraph:P400:end

% mdparagraph:P401:begin
\itempara{4}{\textbf{Adjective models and variants.} FB §§55–60, pp. 38–41; personal substantival adjectives, §181, p. 100. Cook's §57.3 allows internal-vowel variation; the chapter selects the explicitly printed retained-vowel stems consistently. The genitive plural swīðra applies the consonant-cluster simplification of §35.c, p. 25, with the earlier comparative model; it is a declared derivation rather than a claim that Cook prints that whole swīðra table. The distinct instruments of the new adjectives are reference forms. Demonstrative instruments: §§83–84, p. 50. Bare means and mid: §§174–178, pp. 97–99. Cook notes the especially poetic use of bare instruments and the Anglian tendency to put a distinct instrument after mid; the latter is not the prose target here.}
% mdparagraph:P401:end

% mdparagraph:P402:begin
\itempara{5}{\textbf{Verbs and clause constructions.} Finite endings: FB §§93–95, pp. 54–55; dental changes: §§33–35, pp. 24–25. Weak I: §§113,117, pp. 68–69,72; weak II: §§118–119, pp. 73–74; strong I: §102, p. 59. Unnan: §129, p. 77 /\allowbreak{} PDF p. 93; subjunctive commands: §193.a, p. 103. Compare Wright, \emph{Old English Grammar} (1914), §542, p. 284, and Fulk, \emph{An Introductory Grammar of Old English} (2020), §123, pp. 89–90: the absent assigned indicative second singular and the subjunctive command analysis are explicit evidence limits. Cook's label “Imper.” for unne is retained alongside Fulk's explanation.}
% mdparagraph:P402:end

% mdparagraph:P403:begin
Impersonal constructions: FB §§189–191, p. 102, with Fulk §66, p. 42. Present with future reference: FB §188, p. 101. Perception: §169, p. 96, and predicative constructions §173, p. 97. Cook §172.1, p. 97, warns against treating motion/\allowbreak{}rest as an invariable predictor of prepositional case. The bounded on … hand idiom is checked against p. 134, including its static standað in the unassigned continuation. Infinitives: §§198–200, pp. 105–106. The give-to-drink help is adapted from p. 136, editorial lines 5–6: the source is a negative past statement, while III.14.11 requests a positive singular command. The swā lange swā help draws on p. 135, editorial lines 12–14 and note 17; Cook explicitly explains the inasmuch-as force. Only these bounded expressions are supplied, not an unrestricted purpose or temporal-clause system.
% mdparagraph:P403:end

% mdparagraph:P404:begin
\itempara{6}{\textbf{Nominal and adjective lexical evidence.} The following FB vocabulary locators establish the headwords, genders and selected meanings used above. Joined course headwords omit the glossary's segmentation hyphens while preserving letters and quantities. Their paradigms combine these lexical diagnostics with the complete grammatical models cited in notes 3–4.}
% mdparagraph:P404:end

% mdtable:T40:begin
\par\begingroup
\fontsize{10.5}{13.2}\selectfont
\setlength{\LTleft}{0pt}\setlength{\LTright}{\fill}
\begin{longtable}{@{}>{\raggedright\arraybackslash}p{28.73mm}>{\raggedright\arraybackslash}p{26.12mm}>{\raggedright\arraybackslash}p{26.12mm}>{\raggedright\arraybackslash}p{34.6mm}@{}}
\toprule
\textbf{Headword} & \textbf{Printed page} & \textbf{PDF page} & \textbf{Vocabulary entry} \\
\midrule
\endfirsthead
\toprule
\textbf{Headword} & \textbf{Printed page} & \textbf{PDF page} & \textbf{Vocabulary entry} \\
\midrule
\endhead
\midrule
\multicolumn{4}{r}{\footnotesize Continued on the following page.}\\
\endfoot
\bottomrule
\endlastfoot
Crīst & 281 & 297 & s.v. Crīst \\
dryhten & 283 & 299 & s.v. dryhten \\
dōm & 283 & 299 & s.v. dōm \\
rēaf & 309 & 325 & s.v. rēaf \\
cuma & 282 & 298 & s.v. cuma \\
cweartern & 282 & 298 & s.v. cweartern \\
giesthūs & 296 & 312 & s.v. giesthūs \\
fæder & 285 & 301 & s.v. fæder \\
woruld & 323 & 339 & s.v. woruld \\
scēaphierde & 310 & 326 & s.v. scēaphierde \\
scēap & 310 & 326 & s.v. scēap \\
setl & 312 & 328 & s.v. setl \\
hand & 297 & 313 & s.v. hand \\
ðēod & 318 & 334 & s.v. ðēod \\
dēofol & 282 & 298 & s.v. dēofol \\
wǣd & 320 & 336 & s.v. wǣd \\
gāt & 290 & 306 & s.v. gāt \\
nama & 306 & 322 & s.v. nama \\
ęngel & 285 & 301 & s.v. ęngel \\
hungrig & 300 & 316 & s.v. hungrig \\
dēorwierðe & 283 & 299 & s.v. dēorwierðe \\
unrihtwīs & 320 & 336 & s.v. unrihtwīs \\
synfull & 315 & 331 & s.v. synfull \\
untrum & 320 & 336 & s.v. untrum \\
sellic & 312 & 328 & s.v. sellic \\
ēce & 284 & 300 & s.v. ēce \\
cwic & 282 & 298 & s.v. cwic \\
nacod & 306 & 322 & s.v. nacod \\
swīðra & 315 & 331 & s.v. swīð \\
\end{longtable}
\par\endgroup

% mdtable:T40:end
% mdparagraph:P405:begin
\itempara{7}{\textbf{Verbal and constructional lexical evidence.} FB vocabulary: andswarian p. 276, beforan p. 278, fēdan and fela p. 286, forwiernan p. 288, genēosian p. 293, gewītan p. 296, hwonne p. 300, scęncan and scrȳdan p. 311, tōgelǣdan p. 316, wyrcean p. 324. Tōgelǣdan's thing-plus-recipient frame also occurs in Reader II, p. 130, editorial line 5. The already learned cīepan, bycgan, blāwan, hīeran, dōn and tōtwǣman keep their established forms; their new uses are explicitly defined in §§5–8.}
% mdparagraph:P405:end

% mdparagraph:P406:begin
Bosworth–Toller, checked dictionary facsimiles, corroborates the selected lexical frames: \textbf{fédan}, main p. 273b and Supplement p. 207b; \textbf{scrýdan}, main p. 843a; \textbf{scencan}, main p. 828b; \textbf{and-swarian}, main p. 41a; \textbf{for-wiernan}, Supplement p. 258a; \textbf{ge-neósian}, main p. 421a; \textbf{hyngrian}, Supplement p. 586a; \textbf{þyncan}, main pp. 1084b–1085a; \textbf{unnan}, main p. 1122b. In ge-neósian, the cited examples with hine establish an accusative visited person; this does not deny the genitive frame of related nēosian. Cook's genēosian quantity remains the chapter's selected spelling. The scrȳdan garment complement is corroborating reference evidence only; the production frame here requires the accusative person.
% mdparagraph:P406:end

% mdparagraph:P407:begin
\itempara{8}{\textbf{Supplied reading vocabulary.} FB vocabulary entries: bearn p. 277, mægenðrymm and mægenðrymnes p. 304, winstre p. 322, witodlice p. 323, samod p. 310, tōsceadan p. 316, gelōgian p. 293, ætforan p. 275. Their exact local senses and forms are determined by p. 134 and its notes, not by an unrestricted glossary translation. Witodlice is connective here (Cook's Latin autem); the two occurrences of ðonne mean “then”. The source feminine plural ealle and supplied plural passive gegaderode are preserved; ordinary selected ealla remains valid elsewhere. The local helps also supply the divided-in-two phrase, the swā-swā comparison, and the elliptical final hand. The new account and all fourteen teaching examples are original teaching material, not adapted quotations.}
% mdparagraph:P407:end

\endgroup
\end{document}
